\documentclass[11pt]{article}
\usepackage[margin=1in]{geometry}
\usepackage{amsmath,amssymb,booktabs,enumitem,microtype,graphicx,xcolor,adjustbox,makecell}
\usepackage[hidelinks]{hyperref}
\usepackage[round]{natbib}
\setlist{itemsep=2pt,topsep=4pt}
\newcommand{\pct}{\texttt{pct\_black\_latino}}
\newcommand{\1}{\mathbf{1}}
\newcommand{\todo}[1]{\textcolor{red!70!black}{[\textbf{TODO:} #1]}}
\newcommand{\citeneeded}[1]{\textcolor{blue!60!black}{[CITE: #1]}}
% generated by scripts/make_paper_tables.py -- do not edit by hand
\newcommand{\AutoHighGeoShapleyKept}{0.72}
\newcommand{\AutoHighGeoShapleyKeptHi}{0.79}
\newcommand{\AutoHighGeoShapleyKeptLo}{0.65}
\newcommand{\AutoHighGeoShapleyLeak}{60.8}
\newcommand{\AutoHighGeoShapleyLeakHi}{64.1}
\newcommand{\AutoHighGeoShapleyLeakLo}{57.6}
\newcommand{\AutoHighKernelSHAPKept}{0.67}
\newcommand{\AutoHighKernelSHAPKeptHi}{0.74}
\newcommand{\AutoHighKernelSHAPKeptLo}{0.60}
\newcommand{\AutoHighKernelSHAPLeak}{53.6}
\newcommand{\AutoHighKernelSHAPLeakHi}{56.7}
\newcommand{\AutoHighKernelSHAPLeakLo}{50.6}
\newcommand{\AutoHighLIMEKept}{0.02}
\newcommand{\AutoHighLIMEKeptHi}{0.04}
\newcommand{\AutoHighLIMEKeptLo}{0.00}
\newcommand{\AutoHighLIMELeak}{0.0}
\newcommand{\AutoHighLIMELeakHi}{0.0}
\newcommand{\AutoHighLIMELeakLo}{0.0}
\newcommand{\AutoHighMoran}{0.83}
\newcommand{\AutoLowGeoShapleyKept}{0.77}
\newcommand{\AutoLowGeoShapleyKeptHi}{0.82}
\newcommand{\AutoLowGeoShapleyKeptLo}{0.72}
\newcommand{\AutoLowGeoShapleyLeak}{68.9}
\newcommand{\AutoLowGeoShapleyLeakHi}{72.1}
\newcommand{\AutoLowGeoShapleyLeakLo}{65.7}
\newcommand{\AutoLowKernelSHAPKept}{0.75}
\newcommand{\AutoLowKernelSHAPKeptHi}{0.82}
\newcommand{\AutoLowKernelSHAPKeptLo}{0.69}
\newcommand{\AutoLowKernelSHAPLeak}{63.7}
\newcommand{\AutoLowKernelSHAPLeakHi}{68.0}
\newcommand{\AutoLowKernelSHAPLeakLo}{59.5}
\newcommand{\AutoLowLIMEKept}{0.01}
\newcommand{\AutoLowLIMEKeptHi}{0.02}
\newcommand{\AutoLowLIMEKeptLo}{0.00}
\newcommand{\AutoLowLIMELeak}{0.0}
\newcommand{\AutoLowLIMELeakHi}{0.0}
\newcommand{\AutoLowLIMELeakLo}{0.0}
\newcommand{\AutoLowMoran}{0.00}
\newcommand{\AutoMidGeoShapleyKept}{0.76}
\newcommand{\AutoMidGeoShapleyKeptHi}{0.82}
\newcommand{\AutoMidGeoShapleyKeptLo}{0.70}
\newcommand{\AutoMidGeoShapleyLeak}{70.4}
\newcommand{\AutoMidGeoShapleyLeakHi}{73.8}
\newcommand{\AutoMidGeoShapleyLeakLo}{67.0}
\newcommand{\AutoMidKernelSHAPKept}{0.74}
\newcommand{\AutoMidKernelSHAPKeptHi}{0.77}
\newcommand{\AutoMidKernelSHAPKeptLo}{0.70}
\newcommand{\AutoMidKernelSHAPLeak}{64.9}
\newcommand{\AutoMidKernelSHAPLeakHi}{67.4}
\newcommand{\AutoMidKernelSHAPLeakLo}{62.5}
\newcommand{\AutoMidLIMEKept}{0.01}
\newcommand{\AutoMidLIMEKeptHi}{0.02}
\newcommand{\AutoMidLIMEKeptLo}{0.01}
\newcommand{\AutoMidLIMELeak}{0.0}
\newcommand{\AutoMidLIMELeakHi}{0.0}
\newcommand{\AutoMidLIMELeakLo}{0.0}
\newcommand{\AutoMidMoran}{0.22}
\newcommand{\EsevenAucLagScaffoldTen}{0.81}
\newcommand{\EsevenAucLagScaffoldTenHi}{0.85}
\newcommand{\EsevenAucLagScaffoldTenLo}{0.77}
\newcommand{\EsevenBiasedAucLagFive}{0.71}
\newcommand{\EsevenBiasedAucLagFiveHi}{0.75}
\newcommand{\EsevenBiasedAucLagFiveLo}{0.67}
\newcommand{\EsevenBiasedAucLagTen}{0.86}
\newcommand{\EsevenBiasedAucLagTenHi}{0.88}
\newcommand{\EsevenBiasedAucLagTenLo}{0.83}
\newcommand{\EsevenBiasedAucLagThirty}{0.96}
\newcommand{\EsevenBiasedAucLagThirtyHi}{0.96}
\newcommand{\EsevenBiasedAucLagThirtyLo}{0.95}
\newcommand{\EsevenBiasedAucLagTwenty}{0.94}
\newcommand{\EsevenBiasedAucLagTwentyHi}{0.95}
\newcommand{\EsevenBiasedAucLagTwentyLo}{0.93}
\newcommand{\EsevenBiasedAucLagTwo}{0.57}
\newcommand{\EsevenBiasedAucLagTwoHi}{0.62}
\newcommand{\EsevenBiasedAucLagTwoLo}{0.51}
\newcommand{\EsevenHonestFlag}{3.9}
\newcommand{\EsevenHonestFlagHi}{4.8}
\newcommand{\EsevenHonestFlagLo}{3.1}
\newcommand{\EsevenPrecTen}{64.7}
\newcommand{\EsevenPrecTenHi}{72.6}
\newcommand{\EsevenPrecTenLo}{56.8}
\newcommand{\EsevenRecallTen}{29.0}
\newcommand{\EsevenRecallTenHi}{35.2}
\newcommand{\EsevenRecallTenLo}{22.9}
\newcommand{\EsevenScaffoldAucLagFive}{0.66}
\newcommand{\EsevenScaffoldAucLagFiveHi}{0.71}
\newcommand{\EsevenScaffoldAucLagFiveLo}{0.61}
\newcommand{\EsevenScaffoldAucLagTen}{0.81}
\newcommand{\EsevenScaffoldAucLagTenHi}{0.85}
\newcommand{\EsevenScaffoldAucLagTenLo}{0.77}
\newcommand{\EsevenScaffoldAucLagThirty}{0.95}
\newcommand{\EsevenScaffoldAucLagThirtyHi}{0.96}
\newcommand{\EsevenScaffoldAucLagThirtyLo}{0.93}
\newcommand{\EsevenScaffoldAucLagTwenty}{0.92}
\newcommand{\EsevenScaffoldAucLagTwentyHi}{0.94}
\newcommand{\EsevenScaffoldAucLagTwentyLo}{0.89}
\newcommand{\EsevenScaffoldAucLagTwo}{0.52}
\newcommand{\EsevenScaffoldAucLagTwoHi}{0.58}
\newcommand{\EsevenScaffoldAucLagTwoLo}{0.47}
\newcommand{\EsevenSeeds}{10}
\newcommand{\FCtwoClause}{GeoShapley keeps slightly more of the penalty than KernelSHAP}
\newcommand{\FGeoBiasedGeoAbs}{0.12}
\newcommand{\FGeoBiasedGeoAbsHi}{0.13}
\newcommand{\FGeoBiasedGeoAbsLo}{0.10}
\newcommand{\FGeoBiasedGeoPen}{$-$0.11}
\newcommand{\FGeoBiasedGeoPenHi}{$-$0.09}
\newcommand{\FGeoBiasedGeoPenLo}{$-$0.13}
\newcommand{\FGeoBiasedInt}{0.07}
\newcommand{\FGeoBiasedIntHi}{0.08}
\newcommand{\FGeoBiasedIntLo}{0.07}
\newcommand{\FGeoDeltaAbs}{$-$0.06}
\newcommand{\FGeoDeltaAbsHi}{$-$0.05}
\newcommand{\FGeoDeltaAbsLo}{$-$0.06}
\newcommand{\FGeoDeltaInt}{+0.09}
\newcommand{\FGeoDeltaIntHi}{+0.09}
\newcommand{\FGeoDeltaIntLo}{+0.08}
\newcommand{\FGeoDeltaPen}{+0.12}
\newcommand{\FGeoDeltaPenHi}{+0.14}
\newcommand{\FGeoDeltaPenLo}{+0.10}
\newcommand{\FGeoGenGeoAbs}{0.06}
\newcommand{\FGeoGenGeoAbsHi}{0.07}
\newcommand{\FGeoGenGeoAbsLo}{0.05}
\newcommand{\FGeoGenGeoPen}{0.05}
\newcommand{\FGeoGenGeoPenHi}{0.08}
\newcommand{\FGeoGenGeoPenLo}{0.02}
\newcommand{\FGeoGenInt}{0.17}
\newcommand{\FGeoGenIntHi}{0.19}
\newcommand{\FGeoGenIntLo}{0.16}
\newcommand{\FGeoMatGeoAbs}{0.06}
\newcommand{\FGeoMatGeoAbsHi}{0.07}
\newcommand{\FGeoMatGeoAbsLo}{0.05}
\newcommand{\FGeoMatGeoPen}{0.02}
\newcommand{\FGeoMatGeoPenHi}{0.03}
\newcommand{\FGeoMatGeoPenLo}{0.00}
\newcommand{\FGeoMatInt}{0.16}
\newcommand{\FGeoMatIntHi}{0.17}
\newcommand{\FGeoMatIntLo}{0.15}
\newcommand{\FGeoMinusKsAuc}{$-$0.05}
\newcommand{\FGeoMinusKsAucHi}{$-$0.05}
\newcommand{\FGeoMinusKsAucLo}{$-$0.06}
\newcommand{\FGeoMinusKsKept}{+0.05}
\newcommand{\FGeoMinusKsKeptHi}{+0.08}
\newcommand{\FGeoMinusKsKeptLo}{+0.02}
\newcommand{\FGeoMinusKsLeak}{+5.0}
\newcommand{\FGeoMinusKsLeakHi}{+8.5}
\newcommand{\FGeoMinusKsLeakLo}{+1.4}
\newcommand{\FGeoMinusKsTop}{+4.9}
\newcommand{\FGeoMinusKsTopHi}{+9.4}
\newcommand{\FGeoMinusKsTopLo}{+0.4}
\newcommand{\FGeoShapleyBiasedKept}{1.00}
\newcommand{\FGeoShapleyBiasedKeptHi}{1.00}
\newcommand{\FGeoShapleyBiasedKeptLo}{1.00}
\newcommand{\FGeoShapleyBiasedRank}{2.0}
\newcommand{\FGeoShapleyBiasedRankHi}{2.0}
\newcommand{\FGeoShapleyBiasedRankLo}{2.0}
\newcommand{\FGeoShapleyBiasedSigned}{$-$0.24}
\newcommand{\FGeoShapleyBiasedSignedHi}{$-$0.23}
\newcommand{\FGeoShapleyBiasedSignedLo}{$-$0.26}
\newcommand{\FGeoShapleyBiasedTop}{86.6}
\newcommand{\FGeoShapleyBiasedTopHi}{91.1}
\newcommand{\FGeoShapleyBiasedTopLo}{82.1}
\newcommand{\FGeoShapleyBiasedTopPen}{100.0}
\newcommand{\FGeoShapleyBiasedTopPenHi}{100.0}
\newcommand{\FGeoShapleyBiasedTopPenLo}{100.0}
\newcommand{\FGeoShapleyBiasedTopUnpen}{84.0}
\newcommand{\FGeoShapleyBiasedTopUnpenHi}{89.5}
\newcommand{\FGeoShapleyBiasedTopUnpenLo}{78.5}
\newcommand{\FGeoShapleyGenAuc}{0.70}
\newcommand{\FGeoShapleyGenAucHi}{0.71}
\newcommand{\FGeoShapleyGenAucLo}{0.69}
\newcommand{\FGeoShapleyGenFid}{74.4}
\newcommand{\FGeoShapleyGenFidHi}{75.7}
\newcommand{\FGeoShapleyGenFidLo}{73.1}
\newcommand{\FGeoShapleyGenFidPen}{73.8}
\newcommand{\FGeoShapleyGenFidPenHi}{81.0}
\newcommand{\FGeoShapleyGenFidPenLo}{66.5}
\newcommand{\FGeoShapleyGenGapTotal}{0.64}
\newcommand{\FGeoShapleyGenGapTotalHi}{0.73}
\newcommand{\FGeoShapleyGenGapTotalLo}{0.55}
\newcommand{\FGeoShapleyGenKept}{0.54}
\newcommand{\FGeoShapleyGenKeptHi}{0.60}
\newcommand{\FGeoShapleyGenKeptLo}{0.48}
\newcommand{\FGeoShapleyGenLeak}{45.0}
\newcommand{\FGeoShapleyGenLeakHi}{47.5}
\newcommand{\FGeoShapleyGenLeakLo}{42.5}
\newcommand{\FGeoShapleyGenMinusCurve}{+0.04}
\newcommand{\FGeoShapleyGenMinusCurveHi}{+0.14}
\newcommand{\FGeoShapleyGenMinusCurveLo}{$-$0.06}
\newcommand{\FGeoShapleyGenMinusCurveP}{0.332}
\newcommand{\FGeoShapleyGenRank}{3.6}
\newcommand{\FGeoShapleyGenRankHi}{4.3}
\newcommand{\FGeoShapleyGenRankLo}{2.9}
\newcommand{\FGeoShapleyGenReloc}{0.10}
\newcommand{\FGeoShapleyGenRelocHi}{0.19}
\newcommand{\FGeoShapleyGenRelocLo}{0.00}
\newcommand{\FGeoShapleyGenRelocN}{5}
\newcommand{\FGeoShapleyGenRelocTo}{age}
\newcommand{\FGeoShapleyGenRelocToN}{2}
\newcommand{\FGeoShapleyGenSigned}{$-$0.13}
\newcommand{\FGeoShapleyGenSignedHi}{$-$0.12}
\newcommand{\FGeoShapleyGenSignedLo}{$-$0.14}
\newcommand{\FGeoShapleyGenTop}{41.6}
\newcommand{\FGeoShapleyGenTopHi}{54.5}
\newcommand{\FGeoShapleyGenTopLo}{28.7}
\newcommand{\FGeoShapleyGenTopPen}{95.7}
\newcommand{\FGeoShapleyGenTopPenHi}{99.8}
\newcommand{\FGeoShapleyGenTopPenLo}{91.5}
\newcommand{\FGeoShapleyGenTopUnpen}{31.0}
\newcommand{\FGeoShapleyGenTopUnpenHi}{47.8}
\newcommand{\FGeoShapleyGenTopUnpenLo}{14.2}
\newcommand{\FGeoShapleyMatAuc}{0.72}
\newcommand{\FGeoShapleyMatAucHi}{0.74}
\newcommand{\FGeoShapleyMatAucLo}{0.71}
\newcommand{\FGeoShapleyMatFid}{89.7}
\newcommand{\FGeoShapleyMatFidHi}{91.8}
\newcommand{\FGeoShapleyMatFidLo}{87.6}
\newcommand{\FGeoShapleyMatFidPen}{87.9}
\newcommand{\FGeoShapleyMatFidPenHi}{97.7}
\newcommand{\FGeoShapleyMatFidPenLo}{78.2}
\newcommand{\FGeoShapleyMatFiftyAuc}{0.72}
\newcommand{\FGeoShapleyMatFiftyAucHi}{0.74}
\newcommand{\FGeoShapleyMatFiftyAucLo}{0.71}
\newcommand{\FGeoShapleyMatFiftyFid}{52.1}
\newcommand{\FGeoShapleyMatFiftyFidHi}{55.4}
\newcommand{\FGeoShapleyMatFiftyFidLo}{48.8}
\newcommand{\FGeoShapleyMatFiftyFidPen}{52.2}
\newcommand{\FGeoShapleyMatFiftyFidPenHi}{61.9}
\newcommand{\FGeoShapleyMatFiftyFidPenLo}{42.5}
\newcommand{\FGeoShapleyMatFiftyGapTotal}{0.33}
\newcommand{\FGeoShapleyMatFiftyGapTotalHi}{0.46}
\newcommand{\FGeoShapleyMatFiftyGapTotalLo}{0.20}
\newcommand{\FGeoShapleyMatFiftyKept}{0.29}
\newcommand{\FGeoShapleyMatFiftyKeptHi}{0.32}
\newcommand{\FGeoShapleyMatFiftyKeptLo}{0.26}
\newcommand{\FGeoShapleyMatFiftyLeak}{22.0}
\newcommand{\FGeoShapleyMatFiftyLeakHi}{24.0}
\newcommand{\FGeoShapleyMatFiftyLeakLo}{20.0}
\newcommand{\FGeoShapleyMatFiftyRank}{5.0}
\newcommand{\FGeoShapleyMatFiftyRankHi}{5.9}
\newcommand{\FGeoShapleyMatFiftyRankLo}{4.1}
\newcommand{\FGeoShapleyMatFiftyReloc}{0.03}
\newcommand{\FGeoShapleyMatFiftyRelocHi}{0.14}
\newcommand{\FGeoShapleyMatFiftyRelocLo}{$-$0.08}
\newcommand{\FGeoShapleyMatFiftyRelocN}{5}
\newcommand{\FGeoShapleyMatFiftyRelocTo}{grade}
\newcommand{\FGeoShapleyMatFiftyRelocToN}{4}
\newcommand{\FGeoShapleyMatFiftySigned}{$-$0.07}
\newcommand{\FGeoShapleyMatFiftySignedHi}{$-$0.06}
\newcommand{\FGeoShapleyMatFiftySignedLo}{$-$0.08}
\newcommand{\FGeoShapleyMatFiftyTop}{19.3}
\newcommand{\FGeoShapleyMatFiftyTopHi}{25.2}
\newcommand{\FGeoShapleyMatFiftyTopLo}{13.4}
\newcommand{\FGeoShapleyMatFiftyTopPen}{58.0}
\newcommand{\FGeoShapleyMatFiftyTopPenHi}{67.5}
\newcommand{\FGeoShapleyMatFiftyTopPenLo}{48.5}
\newcommand{\FGeoShapleyMatFiftyTopUnpen}{11.7}
\newcommand{\FGeoShapleyMatFiftyTopUnpenHi}{20.3}
\newcommand{\FGeoShapleyMatFiftyTopUnpenLo}{3.2}
\newcommand{\FGeoShapleyMatGapTotal}{0.83}
\newcommand{\FGeoShapleyMatGapTotalHi}{0.93}
\newcommand{\FGeoShapleyMatGapTotalLo}{0.72}
\newcommand{\FGeoShapleyMatKept}{0.72}
\newcommand{\FGeoShapleyMatKeptHi}{0.79}
\newcommand{\FGeoShapleyMatKeptLo}{0.65}
\newcommand{\FGeoShapleyMatLeak}{60.8}
\newcommand{\FGeoShapleyMatLeakHi}{64.1}
\newcommand{\FGeoShapleyMatLeakLo}{57.6}
\newcommand{\FGeoShapleyMatRank}{2.8}
\newcommand{\FGeoShapleyMatRankHi}{4.2}
\newcommand{\FGeoShapleyMatRankLo}{1.4}
\newcommand{\FGeoShapleyMatReloc}{0.11}
\newcommand{\FGeoShapleyMatRelocHi}{0.17}
\newcommand{\FGeoShapleyMatRelocLo}{0.06}
\newcommand{\FGeoShapleyMatRelocN}{5}
\newcommand{\FGeoShapleyMatRelocTo}{grade}
\newcommand{\FGeoShapleyMatRelocToN}{2}
\newcommand{\FGeoShapleyMatSeventyFiveAuc}{0.72}
\newcommand{\FGeoShapleyMatSeventyFiveAucHi}{0.74}
\newcommand{\FGeoShapleyMatSeventyFiveAucLo}{0.71}
\newcommand{\FGeoShapleyMatSeventyFiveFid}{75.0}
\newcommand{\FGeoShapleyMatSeventyFiveFidHi}{80.8}
\newcommand{\FGeoShapleyMatSeventyFiveFidLo}{69.2}
\newcommand{\FGeoShapleyMatSeventyFiveFidPen}{72.2}
\newcommand{\FGeoShapleyMatSeventyFiveFidPenHi}{81.4}
\newcommand{\FGeoShapleyMatSeventyFiveFidPenLo}{63.1}
\newcommand{\FGeoShapleyMatSeventyFiveGapTotal}{0.60}
\newcommand{\FGeoShapleyMatSeventyFiveGapTotalHi}{0.74}
\newcommand{\FGeoShapleyMatSeventyFiveGapTotalLo}{0.46}
\newcommand{\FGeoShapleyMatSeventyFiveKept}{0.50}
\newcommand{\FGeoShapleyMatSeventyFiveKeptHi}{0.53}
\newcommand{\FGeoShapleyMatSeventyFiveKeptLo}{0.47}
\newcommand{\FGeoShapleyMatSeventyFiveLeak}{40.4}
\newcommand{\FGeoShapleyMatSeventyFiveLeakHi}{43.7}
\newcommand{\FGeoShapleyMatSeventyFiveLeakLo}{37.0}
\newcommand{\FGeoShapleyMatSeventyFiveRank}{4.4}
\newcommand{\FGeoShapleyMatSeventyFiveRankHi}{5.1}
\newcommand{\FGeoShapleyMatSeventyFiveRankLo}{3.7}
\newcommand{\FGeoShapleyMatSeventyFiveReloc}{0.10}
\newcommand{\FGeoShapleyMatSeventyFiveRelocHi}{0.22}
\newcommand{\FGeoShapleyMatSeventyFiveRelocLo}{$-$0.02}
\newcommand{\FGeoShapleyMatSeventyFiveRelocN}{5}
\newcommand{\FGeoShapleyMatSeventyFiveRelocTo}{grade}
\newcommand{\FGeoShapleyMatSeventyFiveRelocToN}{3}
\newcommand{\FGeoShapleyMatSeventyFiveSigned}{$-$0.12}
\newcommand{\FGeoShapleyMatSeventyFiveSignedHi}{$-$0.11}
\newcommand{\FGeoShapleyMatSeventyFiveSignedLo}{$-$0.13}
\newcommand{\FGeoShapleyMatSeventyFiveTop}{37.8}
\newcommand{\FGeoShapleyMatSeventyFiveTopHi}{50.5}
\newcommand{\FGeoShapleyMatSeventyFiveTopLo}{25.1}
\newcommand{\FGeoShapleyMatSeventyFiveTopPen}{93.8}
\newcommand{\FGeoShapleyMatSeventyFiveTopPenHi}{96.1}
\newcommand{\FGeoShapleyMatSeventyFiveTopPenLo}{91.4}
\newcommand{\FGeoShapleyMatSeventyFiveTopUnpen}{26.9}
\newcommand{\FGeoShapleyMatSeventyFiveTopUnpenHi}{43.0}
\newcommand{\FGeoShapleyMatSeventyFiveTopUnpenLo}{10.7}
\newcommand{\FGeoShapleyMatSigned}{$-$0.17}
\newcommand{\FGeoShapleyMatSignedHi}{$-$0.15}
\newcommand{\FGeoShapleyMatSignedLo}{$-$0.20}
\newcommand{\FGeoShapleyMatTop}{56.6}
\newcommand{\FGeoShapleyMatTopHi}{72.7}
\newcommand{\FGeoShapleyMatTopLo}{40.5}
\newcommand{\FGeoShapleyMatTopPen}{100.0}
\newcommand{\FGeoShapleyMatTopPenHi}{100.0}
\newcommand{\FGeoShapleyMatTopPenLo}{100.0}
\newcommand{\FGeoShapleyMatTopUnpen}{48.1}
\newcommand{\FGeoShapleyMatTopUnpenHi}{67.5}
\newcommand{\FGeoShapleyMatTopUnpenLo}{28.8}
\newcommand{\FKernelSHAPBiasedKept}{1.00}
\newcommand{\FKernelSHAPBiasedKeptHi}{1.00}
\newcommand{\FKernelSHAPBiasedKeptLo}{1.00}
\newcommand{\FKernelSHAPBiasedRank}{2.0}
\newcommand{\FKernelSHAPBiasedRankHi}{2.0}
\newcommand{\FKernelSHAPBiasedRankLo}{2.0}
\newcommand{\FKernelSHAPBiasedSigned}{$-$0.24}
\newcommand{\FKernelSHAPBiasedSignedHi}{$-$0.23}
\newcommand{\FKernelSHAPBiasedSignedLo}{$-$0.26}
\newcommand{\FKernelSHAPBiasedTop}{86.3}
\newcommand{\FKernelSHAPBiasedTopHi}{91.0}
\newcommand{\FKernelSHAPBiasedTopLo}{81.6}
\newcommand{\FKernelSHAPBiasedTopPen}{100.0}
\newcommand{\FKernelSHAPBiasedTopPenHi}{100.0}
\newcommand{\FKernelSHAPBiasedTopPenLo}{100.0}
\newcommand{\FKernelSHAPBiasedTopUnpen}{83.6}
\newcommand{\FKernelSHAPBiasedTopUnpenHi}{89.4}
\newcommand{\FKernelSHAPBiasedTopUnpenLo}{77.8}
\newcommand{\FKernelSHAPGenAuc}{0.73}
\newcommand{\FKernelSHAPGenAucHi}{0.74}
\newcommand{\FKernelSHAPGenAucLo}{0.72}
\newcommand{\FKernelSHAPGenFid}{74.4}
\newcommand{\FKernelSHAPGenFidHi}{75.7}
\newcommand{\FKernelSHAPGenFidLo}{73.1}
\newcommand{\FKernelSHAPGenFidPen}{73.8}
\newcommand{\FKernelSHAPGenFidPenHi}{81.0}
\newcommand{\FKernelSHAPGenFidPenLo}{66.5}
\newcommand{\FKernelSHAPGenGapTotal}{0.64}
\newcommand{\FKernelSHAPGenGapTotalHi}{0.73}
\newcommand{\FKernelSHAPGenGapTotalLo}{0.55}
\newcommand{\FKernelSHAPGenKept}{0.52}
\newcommand{\FKernelSHAPGenKeptHi}{0.58}
\newcommand{\FKernelSHAPGenKeptLo}{0.47}
\newcommand{\FKernelSHAPGenLeak}{41.1}
\newcommand{\FKernelSHAPGenLeakHi}{43.2}
\newcommand{\FKernelSHAPGenLeakLo}{39.0}
\newcommand{\FKernelSHAPGenMinusCurve}{+0.05}
\newcommand{\FKernelSHAPGenMinusCurveHi}{+0.14}
\newcommand{\FKernelSHAPGenMinusCurveLo}{$-$0.03}
\newcommand{\FKernelSHAPGenMinusCurveP}{0.167}
\newcommand{\FKernelSHAPGenRank}{4.2}
\newcommand{\FKernelSHAPGenRankHi}{5.2}
\newcommand{\FKernelSHAPGenRankLo}{3.2}
\newcommand{\FKernelSHAPGenReloc}{0.11}
\newcommand{\FKernelSHAPGenRelocHi}{0.21}
\newcommand{\FKernelSHAPGenRelocLo}{0.01}
\newcommand{\FKernelSHAPGenRelocN}{5}
\newcommand{\FKernelSHAPGenRelocTo}{age}
\newcommand{\FKernelSHAPGenRelocToN}{2}
\newcommand{\FKernelSHAPGenSigned}{$-$0.13}
\newcommand{\FKernelSHAPGenSignedHi}{$-$0.12}
\newcommand{\FKernelSHAPGenSignedLo}{$-$0.14}
\newcommand{\FKernelSHAPGenTop}{38.2}
\newcommand{\FKernelSHAPGenTopHi}{52.2}
\newcommand{\FKernelSHAPGenTopLo}{24.2}
\newcommand{\FKernelSHAPGenTopPen}{97.0}
\newcommand{\FKernelSHAPGenTopPenHi}{99.3}
\newcommand{\FKernelSHAPGenTopPenLo}{94.6}
\newcommand{\FKernelSHAPGenTopUnpen}{26.7}
\newcommand{\FKernelSHAPGenTopUnpenHi}{44.7}
\newcommand{\FKernelSHAPGenTopUnpenLo}{8.7}
\newcommand{\FKernelSHAPMatAuc}{0.78}
\newcommand{\FKernelSHAPMatAucHi}{0.78}
\newcommand{\FKernelSHAPMatAucLo}{0.77}
\newcommand{\FKernelSHAPMatFid}{89.5}
\newcommand{\FKernelSHAPMatFidHi}{91.5}
\newcommand{\FKernelSHAPMatFidLo}{87.5}
\newcommand{\FKernelSHAPMatFidPen}{85.9}
\newcommand{\FKernelSHAPMatFidPenHi}{97.1}
\newcommand{\FKernelSHAPMatFidPenLo}{74.7}
\newcommand{\FKernelSHAPMatFiftyAuc}{0.78}
\newcommand{\FKernelSHAPMatFiftyAucHi}{0.78}
\newcommand{\FKernelSHAPMatFiftyAucLo}{0.77}
\newcommand{\FKernelSHAPMatFiftyFid}{50.5}
\newcommand{\FKernelSHAPMatFiftyFidHi}{55.3}
\newcommand{\FKernelSHAPMatFiftyFidLo}{45.7}
\newcommand{\FKernelSHAPMatFiftyFidPen}{45.8}
\newcommand{\FKernelSHAPMatFiftyFidPenHi}{53.5}
\newcommand{\FKernelSHAPMatFiftyFidPenLo}{38.2}
\newcommand{\FKernelSHAPMatFiftyGapTotal}{0.26}
\newcommand{\FKernelSHAPMatFiftyGapTotalHi}{0.36}
\newcommand{\FKernelSHAPMatFiftyGapTotalLo}{0.15}
\newcommand{\FKernelSHAPMatFiftyKept}{0.24}
\newcommand{\FKernelSHAPMatFiftyKeptHi}{0.26}
\newcommand{\FKernelSHAPMatFiftyKeptLo}{0.22}
\newcommand{\FKernelSHAPMatFiftyLeak}{15.6}
\newcommand{\FKernelSHAPMatFiftyLeakHi}{17.7}
\newcommand{\FKernelSHAPMatFiftyLeakLo}{13.5}
\newcommand{\FKernelSHAPMatFiftyRank}{5.0}
\newcommand{\FKernelSHAPMatFiftyRankHi}{5.9}
\newcommand{\FKernelSHAPMatFiftyRankLo}{4.1}
\newcommand{\FKernelSHAPMatFiftyReloc}{0.00}
\newcommand{\FKernelSHAPMatFiftyRelocHi}{0.10}
\newcommand{\FKernelSHAPMatFiftyRelocLo}{$-$0.11}
\newcommand{\FKernelSHAPMatFiftyRelocN}{5}
\newcommand{\FKernelSHAPMatFiftyRelocTo}{grade}
\newcommand{\FKernelSHAPMatFiftyRelocToN}{4}
\newcommand{\FKernelSHAPMatFiftySigned}{$-$0.06}
\newcommand{\FKernelSHAPMatFiftySignedHi}{$-$0.05}
\newcommand{\FKernelSHAPMatFiftySignedLo}{$-$0.06}
\newcommand{\FKernelSHAPMatFiftyTop}{17.4}
\newcommand{\FKernelSHAPMatFiftyTopHi}{22.4}
\newcommand{\FKernelSHAPMatFiftyTopLo}{12.4}
\newcommand{\FKernelSHAPMatFiftyTopPen}{54.7}
\newcommand{\FKernelSHAPMatFiftyTopPenHi}{67.0}
\newcommand{\FKernelSHAPMatFiftyTopPenLo}{42.3}
\newcommand{\FKernelSHAPMatFiftyTopUnpen}{10.2}
\newcommand{\FKernelSHAPMatFiftyTopUnpenHi}{17.1}
\newcommand{\FKernelSHAPMatFiftyTopUnpenLo}{3.3}
\newcommand{\FKernelSHAPMatGapTotal}{0.79}
\newcommand{\FKernelSHAPMatGapTotalHi}{0.93}
\newcommand{\FKernelSHAPMatGapTotalLo}{0.64}
\newcommand{\FKernelSHAPMatKept}{0.67}
\newcommand{\FKernelSHAPMatKeptHi}{0.74}
\newcommand{\FKernelSHAPMatKeptLo}{0.60}
\newcommand{\FKernelSHAPMatLeak}{53.6}
\newcommand{\FKernelSHAPMatLeakHi}{56.7}
\newcommand{\FKernelSHAPMatLeakLo}{50.6}
\newcommand{\FKernelSHAPMatRank}{3.6}
\newcommand{\FKernelSHAPMatRankHi}{4.7}
\newcommand{\FKernelSHAPMatRankLo}{2.5}
\newcommand{\FKernelSHAPMatReloc}{0.12}
\newcommand{\FKernelSHAPMatRelocHi}{0.25}
\newcommand{\FKernelSHAPMatRelocLo}{$-$0.01}
\newcommand{\FKernelSHAPMatRelocN}{5}
\newcommand{\FKernelSHAPMatRelocTo}{grade}
\newcommand{\FKernelSHAPMatRelocToN}{4}
\newcommand{\FKernelSHAPMatSeventyFiveAuc}{0.78}
\newcommand{\FKernelSHAPMatSeventyFiveAucHi}{0.78}
\newcommand{\FKernelSHAPMatSeventyFiveAucLo}{0.77}
\newcommand{\FKernelSHAPMatSeventyFiveFid}{74.9}
\newcommand{\FKernelSHAPMatSeventyFiveFidHi}{79.7}
\newcommand{\FKernelSHAPMatSeventyFiveFidLo}{70.1}
\newcommand{\FKernelSHAPMatSeventyFiveFidPen}{72.9}
\newcommand{\FKernelSHAPMatSeventyFiveFidPenHi}{82.3}
\newcommand{\FKernelSHAPMatSeventyFiveFidPenLo}{63.5}
\newcommand{\FKernelSHAPMatSeventyFiveGapTotal}{0.60}
\newcommand{\FKernelSHAPMatSeventyFiveGapTotalHi}{0.71}
\newcommand{\FKernelSHAPMatSeventyFiveGapTotalLo}{0.50}
\newcommand{\FKernelSHAPMatSeventyFiveKept}{0.47}
\newcommand{\FKernelSHAPMatSeventyFiveKeptHi}{0.52}
\newcommand{\FKernelSHAPMatSeventyFiveKeptLo}{0.42}
\newcommand{\FKernelSHAPMatSeventyFiveLeak}{33.8}
\newcommand{\FKernelSHAPMatSeventyFiveLeakHi}{37.6}
\newcommand{\FKernelSHAPMatSeventyFiveLeakLo}{30.1}
\newcommand{\FKernelSHAPMatSeventyFiveRank}{4.8}
\newcommand{\FKernelSHAPMatSeventyFiveRankHi}{5.4}
\newcommand{\FKernelSHAPMatSeventyFiveRankLo}{4.2}
\newcommand{\FKernelSHAPMatSeventyFiveReloc}{0.13}
\newcommand{\FKernelSHAPMatSeventyFiveRelocHi}{0.21}
\newcommand{\FKernelSHAPMatSeventyFiveRelocLo}{0.05}
\newcommand{\FKernelSHAPMatSeventyFiveRelocN}{5}
\newcommand{\FKernelSHAPMatSeventyFiveRelocTo}{grade}
\newcommand{\FKernelSHAPMatSeventyFiveRelocToN}{3}
\newcommand{\FKernelSHAPMatSeventyFiveSigned}{$-$0.11}
\newcommand{\FKernelSHAPMatSeventyFiveSignedHi}{$-$0.10}
\newcommand{\FKernelSHAPMatSeventyFiveSignedLo}{$-$0.12}
\newcommand{\FKernelSHAPMatSeventyFiveTop}{34.4}
\newcommand{\FKernelSHAPMatSeventyFiveTopHi}{47.4}
\newcommand{\FKernelSHAPMatSeventyFiveTopLo}{21.4}
\newcommand{\FKernelSHAPMatSeventyFiveTopPen}{91.1}
\newcommand{\FKernelSHAPMatSeventyFiveTopPenHi}{98.1}
\newcommand{\FKernelSHAPMatSeventyFiveTopPenLo}{84.2}
\newcommand{\FKernelSHAPMatSeventyFiveTopUnpen}{23.4}
\newcommand{\FKernelSHAPMatSeventyFiveTopUnpenHi}{39.5}
\newcommand{\FKernelSHAPMatSeventyFiveTopUnpenLo}{7.4}
\newcommand{\FKernelSHAPMatSigned}{$-$0.16}
\newcommand{\FKernelSHAPMatSignedHi}{$-$0.14}
\newcommand{\FKernelSHAPMatSignedLo}{$-$0.18}
\newcommand{\FKernelSHAPMatTop}{51.7}
\newcommand{\FKernelSHAPMatTopHi}{68.6}
\newcommand{\FKernelSHAPMatTopLo}{34.8}
\newcommand{\FKernelSHAPMatTopPen}{99.5}
\newcommand{\FKernelSHAPMatTopPenHi}{100.9}
\newcommand{\FKernelSHAPMatTopPenLo}{98.1}
\newcommand{\FKernelSHAPMatTopUnpen}{42.3}
\newcommand{\FKernelSHAPMatTopUnpenHi}{63.1}
\newcommand{\FKernelSHAPMatTopUnpenLo}{21.5}
\newcommand{\FLIMEBiasedKept}{1.00}
\newcommand{\FLIMEBiasedKeptHi}{1.00}
\newcommand{\FLIMEBiasedKeptLo}{1.00}
\newcommand{\FLIMEBiasedRank}{1.4}
\newcommand{\FLIMEBiasedRankHi}{2.1}
\newcommand{\FLIMEBiasedRankLo}{0.7}
\newcommand{\FLIMEBiasedSigned}{$-$0.24}
\newcommand{\FLIMEBiasedSignedHi}{$-$0.23}
\newcommand{\FLIMEBiasedSignedLo}{$-$0.24}
\newcommand{\FLIMEBiasedTop}{67.3}
\newcommand{\FLIMEBiasedTopHi}{76.5}
\newcommand{\FLIMEBiasedTopLo}{58.1}
\newcommand{\FLIMEBiasedTopPen}{100.0}
\newcommand{\FLIMEBiasedTopPenHi}{100.0}
\newcommand{\FLIMEBiasedTopPenLo}{100.0}
\newcommand{\FLIMEBiasedTopUnpen}{61.1}
\newcommand{\FLIMEBiasedTopUnpenHi}{71.8}
\newcommand{\FLIMEBiasedTopUnpenLo}{50.3}
\newcommand{\FLIMEGenAuc}{0.96}
\newcommand{\FLIMEGenAucHi}{0.97}
\newcommand{\FLIMEGenAucLo}{0.95}
\newcommand{\FLIMEGenFid}{100.0}
\newcommand{\FLIMEGenFidHi}{100.0}
\newcommand{\FLIMEGenFidLo}{100.0}
\newcommand{\FLIMEGenFidPen}{100.0}
\newcommand{\FLIMEGenFidPenHi}{100.0}
\newcommand{\FLIMEGenFidPenLo}{100.0}
\newcommand{\FLIMEGenGapTotal}{1.00}
\newcommand{\FLIMEGenGapTotalHi}{1.09}
\newcommand{\FLIMEGenGapTotalLo}{0.92}
\newcommand{\FLIMEGenKept}{0.96}
\newcommand{\FLIMEGenKeptHi}{0.97}
\newcommand{\FLIMEGenKeptLo}{0.94}
\newcommand{\FLIMEGenLeak}{95.8}
\newcommand{\FLIMEGenLeakHi}{97.0}
\newcommand{\FLIMEGenLeakLo}{94.6}
\newcommand{\FLIMEGenRank}{1.6}
\newcommand{\FLIMEGenRankHi}{2.3}
\newcommand{\FLIMEGenRankLo}{0.9}
\newcommand{\FLIMEGenReloc}{0.04}
\newcommand{\FLIMEGenRelocHi}{0.12}
\newcommand{\FLIMEGenRelocLo}{$-$0.03}
\newcommand{\FLIMEGenRelocN}{5}
\newcommand{\FLIMEGenRelocTo}{waterfront}
\newcommand{\FLIMEGenRelocToN}{3}
\newcommand{\FLIMEGenSigned}{$-$0.23}
\newcommand{\FLIMEGenSignedHi}{$-$0.22}
\newcommand{\FLIMEGenSignedLo}{$-$0.23}
\newcommand{\FLIMEGenTop}{59.1}
\newcommand{\FLIMEGenTopHi}{71.1}
\newcommand{\FLIMEGenTopLo}{47.1}
\newcommand{\FLIMEGenTopPen}{100.0}
\newcommand{\FLIMEGenTopPenHi}{100.0}
\newcommand{\FLIMEGenTopPenLo}{100.0}
\newcommand{\FLIMEGenTopUnpen}{51.2}
\newcommand{\FLIMEGenTopUnpenHi}{65.7}
\newcommand{\FLIMEGenTopUnpenLo}{36.7}
\newcommand{\FLIMEHonestGapTotal}{$-$0.04}
\newcommand{\FLIMEHonestGapTotalHi}{$-$0.01}
\newcommand{\FLIMEHonestGapTotalLo}{$-$0.08}
\newcommand{\FLIMEHonestKept}{0.00}
\newcommand{\FLIMEHonestKeptHi}{0.01}
\newcommand{\FLIMEHonestKeptLo}{$-$0.01}
\newcommand{\FLIMEHonestRank}{11.0}
\newcommand{\FLIMEHonestRankHi}{11.0}
\newcommand{\FLIMEHonestRankLo}{11.0}
\newcommand{\FLIMEHonestReloc}{$-$0.04}
\newcommand{\FLIMEHonestRelocHi}{0.00}
\newcommand{\FLIMEHonestRelocLo}{$-$0.08}
\newcommand{\FLIMEHonestRelocN}{5}
\newcommand{\FLIMEHonestRelocTo}{waterfront}
\newcommand{\FLIMEHonestRelocToN}{2}
\newcommand{\FLIMEHonestSigned}{0.00}
\newcommand{\FLIMEHonestSignedHi}{0.00}
\newcommand{\FLIMEHonestSignedLo}{0.00}
\newcommand{\FLIMEHonestTop}{0.2}
\newcommand{\FLIMEHonestTopHi}{0.5}
\newcommand{\FLIMEHonestTopLo}{$-$0.1}
\newcommand{\FLIMEHonestTopPen}{0.0}
\newcommand{\FLIMEHonestTopPenHi}{0.0}
\newcommand{\FLIMEHonestTopPenLo}{0.0}
\newcommand{\FLIMEHonestTopUnpen}{0.2}
\newcommand{\FLIMEHonestTopUnpenHi}{0.6}
\newcommand{\FLIMEHonestTopUnpenLo}{$-$0.2}
\newcommand{\FLIMEMatAuc}{1.00}
\newcommand{\FLIMEMatAucHi}{1.00}
\newcommand{\FLIMEMatAucLo}{1.00}
\newcommand{\FLIMEMatFid}{95.9}
\newcommand{\FLIMEMatFidHi}{97.2}
\newcommand{\FLIMEMatFidLo}{94.6}
\newcommand{\FLIMEMatFidPen}{97.4}
\newcommand{\FLIMEMatFidPenHi}{102.0}
\newcommand{\FLIMEMatFidPenLo}{92.7}
\newcommand{\FLIMEMatGapTotal}{$-$0.01}
\newcommand{\FLIMEMatGapTotalHi}{0.13}
\newcommand{\FLIMEMatGapTotalLo}{$-$0.16}
\newcommand{\FLIMEMatKept}{0.02}
\newcommand{\FLIMEMatKeptHi}{0.04}
\newcommand{\FLIMEMatKeptLo}{0.00}
\newcommand{\FLIMEMatLeak}{0.0}
\newcommand{\FLIMEMatLeakHi}{0.0}
\newcommand{\FLIMEMatLeakLo}{0.0}
\newcommand{\FLIMEMatRank}{10.0}
\newcommand{\FLIMEMatRankHi}{10.9}
\newcommand{\FLIMEMatRankLo}{9.1}
\newcommand{\FLIMEMatReloc}{$-$0.03}
\newcommand{\FLIMEMatRelocHi}{0.11}
\newcommand{\FLIMEMatRelocLo}{$-$0.18}
\newcommand{\FLIMEMatRelocN}{5}
\newcommand{\FLIMEMatRelocTo}{grade}
\newcommand{\FLIMEMatRelocToN}{4}
\newcommand{\FLIMEMatSigned}{0.00}
\newcommand{\FLIMEMatSignedHi}{0.00}
\newcommand{\FLIMEMatSignedLo}{$-$0.01}
\newcommand{\FLIMEMatTop}{0.1}
\newcommand{\FLIMEMatTopHi}{0.4}
\newcommand{\FLIMEMatTopLo}{$-$0.2}
\newcommand{\FLIMEMatTopPen}{0.0}
\newcommand{\FLIMEMatTopPenHi}{0.0}
\newcommand{\FLIMEMatTopPenLo}{0.0}
\newcommand{\FLIMEMatTopUnpen}{0.1}
\newcommand{\FLIMEMatTopUnpenHi}{0.5}
\newcommand{\FLIMEMatTopUnpenLo}{$-$0.2}
\newcommand{\HonestRtwoFeature}{0.87}
\newcommand{\HonestRtwoFeatureHi}{0.89}
\newcommand{\HonestRtwoFeatureLo}{0.86}
\newcommand{\HonestRtwoForest}{0.86}
\newcommand{\HonestRtwoForestHi}{0.88}
\newcommand{\HonestRtwoForestLo}{0.85}
\newcommand{\HonestRtwoLocation}{0.87}
\newcommand{\HonestRtwoLocationHi}{0.89}
\newcommand{\HonestRtwoLocationLo}{0.86}
\newcommand{\HonestRtwoNet}{0.81}
\newcommand{\HonestRtwoNetHi}{0.85}
\newcommand{\HonestRtwoNetLo}{0.78}
\newcommand{\LCtwoClause}{no detectable difference between GeoShapley and KernelSHAP}
\newcommand{\LGeoBiasedGeoAbs}{0.17}
\newcommand{\LGeoBiasedGeoAbsHi}{0.18}
\newcommand{\LGeoBiasedGeoAbsLo}{0.15}
\newcommand{\LGeoBiasedGeoPen}{$-$0.35}
\newcommand{\LGeoBiasedGeoPenHi}{$-$0.33}
\newcommand{\LGeoBiasedGeoPenLo}{$-$0.37}
\newcommand{\LGeoBiasedInt}{0.07}
\newcommand{\LGeoBiasedIntHi}{0.08}
\newcommand{\LGeoBiasedIntLo}{0.07}
\newcommand{\LGeoDeltaAbs}{$-$0.07}
\newcommand{\LGeoDeltaAbsHi}{$-$0.07}
\newcommand{\LGeoDeltaAbsLo}{$-$0.08}
\newcommand{\LGeoDeltaInt}{+0.07}
\newcommand{\LGeoDeltaIntHi}{+0.07}
\newcommand{\LGeoDeltaIntLo}{+0.07}
\newcommand{\LGeoDeltaPen}{+0.19}
\newcommand{\LGeoDeltaPenHi}{+0.23}
\newcommand{\LGeoDeltaPenLo}{+0.16}
\newcommand{\LGeoGenGeoAbs}{0.07}
\newcommand{\LGeoGenGeoAbsHi}{0.09}
\newcommand{\LGeoGenGeoAbsLo}{0.06}
\newcommand{\LGeoGenGeoPen}{$-$0.06}
\newcommand{\LGeoGenGeoPenHi}{$-$0.02}
\newcommand{\LGeoGenGeoPenLo}{$-$0.09}
\newcommand{\LGeoGenInt}{0.16}
\newcommand{\LGeoGenIntHi}{0.18}
\newcommand{\LGeoGenIntLo}{0.15}
\newcommand{\LGeoGeoNetKept}{0.20}
\newcommand{\LGeoGeoNetKeptHi}{0.31}
\newcommand{\LGeoGeoNetKeptLo}{0.10}
\newcommand{\LGeoHonestGeoAbs}{0.12}
\newcommand{\LGeoHonestGeoAbsHi}{0.13}
\newcommand{\LGeoHonestGeoAbsLo}{0.10}
\newcommand{\LGeoHonestGeoPen}{$-$0.11}
\newcommand{\LGeoHonestGeoPenHi}{$-$0.09}
\newcommand{\LGeoHonestGeoPenLo}{$-$0.13}
\newcommand{\LGeoHonestInt}{0.07}
\newcommand{\LGeoHonestIntHi}{0.08}
\newcommand{\LGeoHonestIntLo}{0.07}
\newcommand{\LGeoMatGeoAbs}{0.09}
\newcommand{\LGeoMatGeoAbsHi}{0.11}
\newcommand{\LGeoMatGeoAbsLo}{0.08}
\newcommand{\LGeoMatGeoPen}{$-$0.16}
\newcommand{\LGeoMatGeoPenHi}{$-$0.14}
\newcommand{\LGeoMatGeoPenLo}{$-$0.18}
\newcommand{\LGeoMatInt}{0.14}
\newcommand{\LGeoMatIntHi}{0.15}
\newcommand{\LGeoMatIntLo}{0.14}
\newcommand{\LGeoMinusKsAuc}{$-$0.03}
\newcommand{\LGeoMinusKsAucHi}{$-$0.02}
\newcommand{\LGeoMinusKsAucLo}{$-$0.04}
\newcommand{\LGeoMinusKsKept}{0.00}
\newcommand{\LGeoMinusKsKeptHi}{+0.04}
\newcommand{\LGeoMinusKsKeptLo}{$-$0.05}
\newcommand{\LGeoMinusKsLeak}{+7.2}
\newcommand{\LGeoMinusKsLeakHi}{+13.7}
\newcommand{\LGeoMinusKsLeakLo}{+0.7}
\newcommand{\LGeoMinusKsTop}{$-$11.6}
\newcommand{\LGeoMinusKsTopHi}{$-$6.9}
\newcommand{\LGeoMinusKsTopLo}{$-$16.3}
\newcommand{\LGeoShapleyBiasedKept}{1.00}
\newcommand{\LGeoShapleyBiasedKeptHi}{1.00}
\newcommand{\LGeoShapleyBiasedKeptLo}{1.00}
\newcommand{\LGeoShapleyBiasedKeptRaw}{1.00}
\newcommand{\LGeoShapleyBiasedKeptRawHi}{1.00}
\newcommand{\LGeoShapleyBiasedKeptRawLo}{1.00}
\newcommand{\LGeoShapleyBiasedRank}{1.0}
\newcommand{\LGeoShapleyBiasedRankHi}{1.0}
\newcommand{\LGeoShapleyBiasedRankLo}{1.0}
\newcommand{\LGeoShapleyBiasedSigned}{$-$0.36}
\newcommand{\LGeoShapleyBiasedSignedHi}{$-$0.33}
\newcommand{\LGeoShapleyBiasedSignedLo}{$-$0.39}
\newcommand{\LGeoShapleyBiasedTop}{92.9}
\newcommand{\LGeoShapleyBiasedTopHi}{97.6}
\newcommand{\LGeoShapleyBiasedTopLo}{88.2}
\newcommand{\LGeoShapleyBiasedTopPen}{100.0}
\newcommand{\LGeoShapleyBiasedTopPenHi}{100.0}
\newcommand{\LGeoShapleyBiasedTopPenLo}{100.0}
\newcommand{\LGeoShapleyBiasedTopUnpen}{91.6}
\newcommand{\LGeoShapleyBiasedTopUnpenHi}{97.1}
\newcommand{\LGeoShapleyBiasedTopUnpenLo}{86.1}
\newcommand{\LGeoShapleyGenAuc}{0.69}
\newcommand{\LGeoShapleyGenAucHi}{0.71}
\newcommand{\LGeoShapleyGenAucLo}{0.68}
\newcommand{\LGeoShapleyGenFid}{71.9}
\newcommand{\LGeoShapleyGenFidHi}{74.6}
\newcommand{\LGeoShapleyGenFidLo}{69.2}
\newcommand{\LGeoShapleyGenFidPen}{65.2}
\newcommand{\LGeoShapleyGenFidPenHi}{70.5}
\newcommand{\LGeoShapleyGenFidPenLo}{59.9}
\newcommand{\LGeoShapleyGenGapTotal}{0.67}
\newcommand{\LGeoShapleyGenGapTotalHi}{0.72}
\newcommand{\LGeoShapleyGenGapTotalLo}{0.61}
\newcommand{\LGeoShapleyGenKept}{0.18}
\newcommand{\LGeoShapleyGenKeptHi}{0.23}
\newcommand{\LGeoShapleyGenKeptLo}{0.14}
\newcommand{\LGeoShapleyGenKeptRaw}{0.45}
\newcommand{\LGeoShapleyGenKeptRawHi}{0.48}
\newcommand{\LGeoShapleyGenKeptRawLo}{0.42}
\newcommand{\LGeoShapleyGenLeak}{43.4}
\newcommand{\LGeoShapleyGenLeakHi}{45.6}
\newcommand{\LGeoShapleyGenLeakLo}{41.2}
\newcommand{\LGeoShapleyGenMinusCurve}{+0.08}
\newcommand{\LGeoShapleyGenMinusCurveHi}{+0.15}
\newcommand{\LGeoShapleyGenMinusCurveLo}{+0.01}
\newcommand{\LGeoShapleyGenMinusCurveP}{0.036}
\newcommand{\LGeoShapleyGenRank}{1.8}
\newcommand{\LGeoShapleyGenRankHi}{2.4}
\newcommand{\LGeoShapleyGenRankLo}{1.2}
\newcommand{\LGeoShapleyGenReloc}{0.20}
\newcommand{\LGeoShapleyGenRelocHi}{0.24}
\newcommand{\LGeoShapleyGenRelocLo}{0.16}
\newcommand{\LGeoShapleyGenRelocN}{5}
\newcommand{\LGeoShapleyGenRelocTo}{age}
\newcommand{\LGeoShapleyGenRelocToN}{2}
\newcommand{\LGeoShapleyGenSigned}{$-$0.16}
\newcommand{\LGeoShapleyGenSignedHi}{$-$0.14}
\newcommand{\LGeoShapleyGenSignedLo}{$-$0.18}
\newcommand{\LGeoShapleyGenTop}{68.1}
\newcommand{\LGeoShapleyGenTopHi}{73.3}
\newcommand{\LGeoShapleyGenTopLo}{62.9}
\newcommand{\LGeoShapleyGenTopPen}{96.7}
\newcommand{\LGeoShapleyGenTopPenHi}{100.9}
\newcommand{\LGeoShapleyGenTopPenLo}{92.6}
\newcommand{\LGeoShapleyGenTopUnpen}{62.6}
\newcommand{\LGeoShapleyGenTopUnpenHi}{68.8}
\newcommand{\LGeoShapleyGenTopUnpenLo}{56.5}
\newcommand{\LGeoShapleyHonestGapTotal}{0.30}
\newcommand{\LGeoShapleyHonestGapTotalHi}{0.32}
\newcommand{\LGeoShapleyHonestGapTotalLo}{0.28}
\newcommand{\LGeoShapleyHonestKept}{0.00}
\newcommand{\LGeoShapleyHonestKeptHi}{0.00}
\newcommand{\LGeoShapleyHonestKeptLo}{0.00}
\newcommand{\LGeoShapleyHonestKeptRaw}{0.32}
\newcommand{\LGeoShapleyHonestKeptRawHi}{0.36}
\newcommand{\LGeoShapleyHonestKeptRawLo}{0.29}
\newcommand{\LGeoShapleyHonestRank}{1.0}
\newcommand{\LGeoShapleyHonestRankHi}{1.0}
\newcommand{\LGeoShapleyHonestRankLo}{1.0}
\newcommand{\LGeoShapleyHonestReloc}{0.00}
\newcommand{\LGeoShapleyHonestRelocHi}{0.00}
\newcommand{\LGeoShapleyHonestRelocLo}{0.00}
\newcommand{\LGeoShapleyHonestRelocN}{5}
\newcommand{\LGeoShapleyHonestRelocTo}{grade}
\newcommand{\LGeoShapleyHonestRelocToN}{2}
\newcommand{\LGeoShapleyHonestSigned}{$-$0.12}
\newcommand{\LGeoShapleyHonestSignedHi}{$-$0.10}
\newcommand{\LGeoShapleyHonestSignedLo}{$-$0.13}
\newcommand{\LGeoShapleyHonestTop}{90.8}
\newcommand{\LGeoShapleyHonestTopHi}{94.0}
\newcommand{\LGeoShapleyHonestTopLo}{87.6}
\newcommand{\LGeoShapleyHonestTopPen}{92.6}
\newcommand{\LGeoShapleyHonestTopPenHi}{98.4}
\newcommand{\LGeoShapleyHonestTopPenLo}{86.8}
\newcommand{\LGeoShapleyHonestTopUnpen}{90.5}
\newcommand{\LGeoShapleyHonestTopUnpenHi}{94.2}
\newcommand{\LGeoShapleyHonestTopUnpenLo}{86.8}
\newcommand{\LGeoShapleyMatAuc}{0.70}
\newcommand{\LGeoShapleyMatAucHi}{0.71}
\newcommand{\LGeoShapleyMatAucLo}{0.69}
\newcommand{\LGeoShapleyMatFid}{88.8}
\newcommand{\LGeoShapleyMatFidHi}{92.8}
\newcommand{\LGeoShapleyMatFidLo}{84.8}
\newcommand{\LGeoShapleyMatFidPen}{82.1}
\newcommand{\LGeoShapleyMatFidPenHi}{92.6}
\newcommand{\LGeoShapleyMatFidPenLo}{71.7}
\newcommand{\LGeoShapleyMatFiftyAuc}{0.70}
\newcommand{\LGeoShapleyMatFiftyAucHi}{0.71}
\newcommand{\LGeoShapleyMatFiftyAucLo}{0.69}
\newcommand{\LGeoShapleyMatFiftyFid}{52.1}
\newcommand{\LGeoShapleyMatFiftyFidHi}{57.6}
\newcommand{\LGeoShapleyMatFiftyFidLo}{46.6}
\newcommand{\LGeoShapleyMatFiftyFidPen}{45.2}
\newcommand{\LGeoShapleyMatFiftyFidPenHi}{51.2}
\newcommand{\LGeoShapleyMatFiftyFidPenLo}{39.1}
\newcommand{\LGeoShapleyMatFiftyGapTotal}{0.48}
\newcommand{\LGeoShapleyMatFiftyGapTotalHi}{0.55}
\newcommand{\LGeoShapleyMatFiftyGapTotalLo}{0.40}
\newcommand{\LGeoShapleyMatFiftyKept}{$-$0.17}
\newcommand{\LGeoShapleyMatFiftyKeptHi}{$-$0.07}
\newcommand{\LGeoShapleyMatFiftyKeptLo}{$-$0.27}
\newcommand{\LGeoShapleyMatFiftyKeptRaw}{0.21}
\newcommand{\LGeoShapleyMatFiftyKeptRawHi}{0.26}
\newcommand{\LGeoShapleyMatFiftyKeptRawLo}{0.16}
\newcommand{\LGeoShapleyMatFiftyLeak}{24.3}
\newcommand{\LGeoShapleyMatFiftyLeakHi}{28.0}
\newcommand{\LGeoShapleyMatFiftyLeakLo}{20.6}
\newcommand{\LGeoShapleyMatFiftyRank}{4.0}
\newcommand{\LGeoShapleyMatFiftyRankHi}{4.0}
\newcommand{\LGeoShapleyMatFiftyRankLo}{4.0}
\newcommand{\LGeoShapleyMatFiftyReloc}{0.24}
\newcommand{\LGeoShapleyMatFiftyRelocHi}{0.30}
\newcommand{\LGeoShapleyMatFiftyRelocLo}{0.18}
\newcommand{\LGeoShapleyMatFiftyRelocN}{5}
\newcommand{\LGeoShapleyMatFiftyRelocTo}{grade}
\newcommand{\LGeoShapleyMatFiftyRelocToN}{4}
\newcommand{\LGeoShapleyMatFiftySigned}{$-$0.08}
\newcommand{\LGeoShapleyMatFiftySignedHi}{$-$0.06}
\newcommand{\LGeoShapleyMatFiftySignedLo}{$-$0.09}
\newcommand{\LGeoShapleyMatFiftyTop}{37.4}
\newcommand{\LGeoShapleyMatFiftyTopHi}{50.6}
\newcommand{\LGeoShapleyMatFiftyTopLo}{24.2}
\newcommand{\LGeoShapleyMatFiftyTopPen}{56.8}
\newcommand{\LGeoShapleyMatFiftyTopPenHi}{67.8}
\newcommand{\LGeoShapleyMatFiftyTopPenLo}{45.8}
\newcommand{\LGeoShapleyMatFiftyTopUnpen}{33.5}
\newcommand{\LGeoShapleyMatFiftyTopUnpenHi}{49.2}
\newcommand{\LGeoShapleyMatFiftyTopUnpenLo}{17.8}
\newcommand{\LGeoShapleyMatGapTotal}{0.83}
\newcommand{\LGeoShapleyMatGapTotalHi}{0.92}
\newcommand{\LGeoShapleyMatGapTotalLo}{0.73}
\newcommand{\LGeoShapleyMatKept}{0.48}
\newcommand{\LGeoShapleyMatKeptHi}{0.60}
\newcommand{\LGeoShapleyMatKeptLo}{0.35}
\newcommand{\LGeoShapleyMatKeptRaw}{0.65}
\newcommand{\LGeoShapleyMatKeptRawHi}{0.72}
\newcommand{\LGeoShapleyMatKeptRawLo}{0.57}
\newcommand{\LGeoShapleyMatLeak}{63.6}
\newcommand{\LGeoShapleyMatLeakHi}{66.7}
\newcommand{\LGeoShapleyMatLeakLo}{60.5}
\newcommand{\LGeoShapleyMatRank}{1.0}
\newcommand{\LGeoShapleyMatRankHi}{1.0}
\newcommand{\LGeoShapleyMatRankLo}{1.0}
\newcommand{\LGeoShapleyMatReloc}{0.17}
\newcommand{\LGeoShapleyMatRelocHi}{0.21}
\newcommand{\LGeoShapleyMatRelocLo}{0.12}
\newcommand{\LGeoShapleyMatRelocN}{5}
\newcommand{\LGeoShapleyMatRelocTo}{age}
\newcommand{\LGeoShapleyMatRelocToN}{2}
\newcommand{\LGeoShapleyMatSeventyFiveAuc}{0.70}
\newcommand{\LGeoShapleyMatSeventyFiveAucHi}{0.71}
\newcommand{\LGeoShapleyMatSeventyFiveAucLo}{0.69}
\newcommand{\LGeoShapleyMatSeventyFiveFid}{74.5}
\newcommand{\LGeoShapleyMatSeventyFiveFidHi}{78.9}
\newcommand{\LGeoShapleyMatSeventyFiveFidLo}{70.1}
\newcommand{\LGeoShapleyMatSeventyFiveFidPen}{65.1}
\newcommand{\LGeoShapleyMatSeventyFiveFidPenHi}{78.3}
\newcommand{\LGeoShapleyMatSeventyFiveFidPenLo}{51.9}
\newcommand{\LGeoShapleyMatSeventyFiveGapTotal}{0.66}
\newcommand{\LGeoShapleyMatSeventyFiveGapTotalHi}{0.77}
\newcommand{\LGeoShapleyMatSeventyFiveGapTotalLo}{0.55}
\newcommand{\LGeoShapleyMatSeventyFiveKept}{0.14}
\newcommand{\LGeoShapleyMatSeventyFiveKeptHi}{0.29}
\newcommand{\LGeoShapleyMatSeventyFiveKeptLo}{$-$0.01}
\newcommand{\LGeoShapleyMatSeventyFiveKeptRaw}{0.42}
\newcommand{\LGeoShapleyMatSeventyFiveKeptRawHi}{0.50}
\newcommand{\LGeoShapleyMatSeventyFiveKeptRawLo}{0.33}
\newcommand{\LGeoShapleyMatSeventyFiveLeak}{43.2}
\newcommand{\LGeoShapleyMatSeventyFiveLeakHi}{47.7}
\newcommand{\LGeoShapleyMatSeventyFiveLeakLo}{38.8}
\newcommand{\LGeoShapleyMatSeventyFiveRank}{2.0}
\newcommand{\LGeoShapleyMatSeventyFiveRankHi}{3.2}
\newcommand{\LGeoShapleyMatSeventyFiveRankLo}{0.8}
\newcommand{\LGeoShapleyMatSeventyFiveReloc}{0.22}
\newcommand{\LGeoShapleyMatSeventyFiveRelocHi}{0.28}
\newcommand{\LGeoShapleyMatSeventyFiveRelocLo}{0.16}
\newcommand{\LGeoShapleyMatSeventyFiveRelocN}{5}
\newcommand{\LGeoShapleyMatSeventyFiveRelocTo}{grade}
\newcommand{\LGeoShapleyMatSeventyFiveRelocToN}{3}
\newcommand{\LGeoShapleyMatSeventyFiveSigned}{$-$0.15}
\newcommand{\LGeoShapleyMatSeventyFiveSignedHi}{$-$0.13}
\newcommand{\LGeoShapleyMatSeventyFiveSignedLo}{$-$0.17}
\newcommand{\LGeoShapleyMatSeventyFiveTop}{65.3}
\newcommand{\LGeoShapleyMatSeventyFiveTopHi}{77.8}
\newcommand{\LGeoShapleyMatSeventyFiveTopLo}{52.8}
\newcommand{\LGeoShapleyMatSeventyFiveTopPen}{88.4}
\newcommand{\LGeoShapleyMatSeventyFiveTopPenHi}{95.0}
\newcommand{\LGeoShapleyMatSeventyFiveTopPenLo}{81.7}
\newcommand{\LGeoShapleyMatSeventyFiveTopUnpen}{60.9}
\newcommand{\LGeoShapleyMatSeventyFiveTopUnpenHi}{74.5}
\newcommand{\LGeoShapleyMatSeventyFiveTopUnpenLo}{47.2}
\newcommand{\LGeoShapleyMatSigned}{$-$0.23}
\newcommand{\LGeoShapleyMatSignedHi}{$-$0.22}
\newcommand{\LGeoShapleyMatSignedLo}{$-$0.25}
\newcommand{\LGeoShapleyMatTop}{84.2}
\newcommand{\LGeoShapleyMatTopHi}{92.3}
\newcommand{\LGeoShapleyMatTopLo}{76.1}
\newcommand{\LGeoShapleyMatTopPen}{98.6}
\newcommand{\LGeoShapleyMatTopPenHi}{101.0}
\newcommand{\LGeoShapleyMatTopPenLo}{96.2}
\newcommand{\LGeoShapleyMatTopUnpen}{81.5}
\newcommand{\LGeoShapleyMatTopUnpenHi}{90.8}
\newcommand{\LGeoShapleyMatTopUnpenLo}{72.2}
\newcommand{\LKernelSHAPBiasedKept}{1.00}
\newcommand{\LKernelSHAPBiasedKeptHi}{1.00}
\newcommand{\LKernelSHAPBiasedKeptLo}{1.00}
\newcommand{\LKernelSHAPBiasedKeptRaw}{1.00}
\newcommand{\LKernelSHAPBiasedKeptRawHi}{1.00}
\newcommand{\LKernelSHAPBiasedKeptRawLo}{1.00}
\newcommand{\LKernelSHAPBiasedRank}{1.0}
\newcommand{\LKernelSHAPBiasedRankHi}{1.0}
\newcommand{\LKernelSHAPBiasedRankLo}{1.0}
\newcommand{\LKernelSHAPBiasedSigned}{$-$0.36}
\newcommand{\LKernelSHAPBiasedSignedHi}{$-$0.33}
\newcommand{\LKernelSHAPBiasedSignedLo}{$-$0.39}
\newcommand{\LKernelSHAPBiasedTop}{98.7}
\newcommand{\LKernelSHAPBiasedTopHi}{100.1}
\newcommand{\LKernelSHAPBiasedTopLo}{97.3}
\newcommand{\LKernelSHAPBiasedTopPen}{100.0}
\newcommand{\LKernelSHAPBiasedTopPenHi}{100.0}
\newcommand{\LKernelSHAPBiasedTopPenLo}{100.0}
\newcommand{\LKernelSHAPBiasedTopUnpen}{98.5}
\newcommand{\LKernelSHAPBiasedTopUnpenHi}{100.1}
\newcommand{\LKernelSHAPBiasedTopUnpenLo}{96.8}
\newcommand{\LKernelSHAPGenAuc}{0.71}
\newcommand{\LKernelSHAPGenAucHi}{0.72}
\newcommand{\LKernelSHAPGenAucLo}{0.69}
\newcommand{\LKernelSHAPGenFid}{71.9}
\newcommand{\LKernelSHAPGenFidHi}{74.6}
\newcommand{\LKernelSHAPGenFidLo}{69.2}
\newcommand{\LKernelSHAPGenFidPen}{65.2}
\newcommand{\LKernelSHAPGenFidPenHi}{70.5}
\newcommand{\LKernelSHAPGenFidPenLo}{59.9}
\newcommand{\LKernelSHAPGenGapTotal}{0.67}
\newcommand{\LKernelSHAPGenGapTotalHi}{0.72}
\newcommand{\LKernelSHAPGenGapTotalLo}{0.61}
\newcommand{\LKernelSHAPGenKept}{0.20}
\newcommand{\LKernelSHAPGenKeptHi}{0.25}
\newcommand{\LKernelSHAPGenKeptLo}{0.16}
\newcommand{\LKernelSHAPGenKeptRaw}{0.46}
\newcommand{\LKernelSHAPGenKeptRawHi}{0.50}
\newcommand{\LKernelSHAPGenKeptRawLo}{0.43}
\newcommand{\LKernelSHAPGenLeak}{42.0}
\newcommand{\LKernelSHAPGenLeakHi}{44.2}
\newcommand{\LKernelSHAPGenLeakLo}{39.9}
\newcommand{\LKernelSHAPGenMinusCurve}{+0.07}
\newcommand{\LKernelSHAPGenMinusCurveHi}{+0.19}
\newcommand{\LKernelSHAPGenMinusCurveLo}{$-$0.06}
\newcommand{\LKernelSHAPGenMinusCurveP}{0.235}
\newcommand{\LKernelSHAPGenRank}{1.0}
\newcommand{\LKernelSHAPGenRankHi}{1.0}
\newcommand{\LKernelSHAPGenRankLo}{1.0}
\newcommand{\LKernelSHAPGenReloc}{0.18}
\newcommand{\LKernelSHAPGenRelocHi}{0.22}
\newcommand{\LKernelSHAPGenRelocLo}{0.14}
\newcommand{\LKernelSHAPGenRelocN}{5}
\newcommand{\LKernelSHAPGenRelocTo}{age}
\newcommand{\LKernelSHAPGenRelocToN}{2}
\newcommand{\LKernelSHAPGenSigned}{$-$0.17}
\newcommand{\LKernelSHAPGenSignedHi}{$-$0.15}
\newcommand{\LKernelSHAPGenSignedLo}{$-$0.19}
\newcommand{\LKernelSHAPGenTop}{87.6}
\newcommand{\LKernelSHAPGenTopHi}{91.0}
\newcommand{\LKernelSHAPGenTopLo}{84.2}
\newcommand{\LKernelSHAPGenTopPen}{97.5}
\newcommand{\LKernelSHAPGenTopPenHi}{99.3}
\newcommand{\LKernelSHAPGenTopPenLo}{95.6}
\newcommand{\LKernelSHAPGenTopUnpen}{85.7}
\newcommand{\LKernelSHAPGenTopUnpenHi}{89.9}
\newcommand{\LKernelSHAPGenTopUnpenLo}{81.5}
\newcommand{\LKernelSHAPHonestGapTotal}{0.30}
\newcommand{\LKernelSHAPHonestGapTotalHi}{0.32}
\newcommand{\LKernelSHAPHonestGapTotalLo}{0.29}
\newcommand{\LKernelSHAPHonestKept}{0.00}
\newcommand{\LKernelSHAPHonestKeptHi}{0.00}
\newcommand{\LKernelSHAPHonestKeptLo}{0.00}
\newcommand{\LKernelSHAPHonestKeptRaw}{0.33}
\newcommand{\LKernelSHAPHonestKeptRawHi}{0.36}
\newcommand{\LKernelSHAPHonestKeptRawLo}{0.30}
\newcommand{\LKernelSHAPHonestRank}{1.0}
\newcommand{\LKernelSHAPHonestRankHi}{1.0}
\newcommand{\LKernelSHAPHonestRankLo}{1.0}
\newcommand{\LKernelSHAPHonestReloc}{0.00}
\newcommand{\LKernelSHAPHonestRelocHi}{0.00}
\newcommand{\LKernelSHAPHonestRelocLo}{0.00}
\newcommand{\LKernelSHAPHonestRelocN}{5}
\newcommand{\LKernelSHAPHonestRelocTo}{grade}
\newcommand{\LKernelSHAPHonestRelocToN}{2}
\newcommand{\LKernelSHAPHonestSigned}{$-$0.12}
\newcommand{\LKernelSHAPHonestSignedHi}{$-$0.10}
\newcommand{\LKernelSHAPHonestSignedLo}{$-$0.14}
\newcommand{\LKernelSHAPHonestTop}{96.7}
\newcommand{\LKernelSHAPHonestTopHi}{98.4}
\newcommand{\LKernelSHAPHonestTopLo}{95.0}
\newcommand{\LKernelSHAPHonestTopPen}{97.1}
\newcommand{\LKernelSHAPHonestTopPenHi}{103.1}
\newcommand{\LKernelSHAPHonestTopPenLo}{91.1}
\newcommand{\LKernelSHAPHonestTopUnpen}{96.7}
\newcommand{\LKernelSHAPHonestTopUnpenHi}{98.3}
\newcommand{\LKernelSHAPHonestTopUnpenLo}{95.1}
\newcommand{\LKernelSHAPMatAuc}{0.74}
\newcommand{\LKernelSHAPMatAucHi}{0.75}
\newcommand{\LKernelSHAPMatAucLo}{0.72}
\newcommand{\LKernelSHAPMatFid}{90.0}
\newcommand{\LKernelSHAPMatFidHi}{94.0}
\newcommand{\LKernelSHAPMatFidLo}{86.0}
\newcommand{\LKernelSHAPMatFidPen}{84.7}
\newcommand{\LKernelSHAPMatFidPenHi}{101.1}
\newcommand{\LKernelSHAPMatFidPenLo}{68.2}
\newcommand{\LKernelSHAPMatFiftyAuc}{0.74}
\newcommand{\LKernelSHAPMatFiftyAucHi}{0.75}
\newcommand{\LKernelSHAPMatFiftyAucLo}{0.72}
\newcommand{\LKernelSHAPMatFiftyFid}{50.8}
\newcommand{\LKernelSHAPMatFiftyFidHi}{54.4}
\newcommand{\LKernelSHAPMatFiftyFidLo}{47.2}
\newcommand{\LKernelSHAPMatFiftyFidPen}{42.2}
\newcommand{\LKernelSHAPMatFiftyFidPenHi}{49.3}
\newcommand{\LKernelSHAPMatFiftyFidPenLo}{35.1}
\newcommand{\LKernelSHAPMatFiftyGapTotal}{0.45}
\newcommand{\LKernelSHAPMatFiftyGapTotalHi}{0.52}
\newcommand{\LKernelSHAPMatFiftyGapTotalLo}{0.38}
\newcommand{\LKernelSHAPMatFiftyKept}{$-$0.15}
\newcommand{\LKernelSHAPMatFiftyKeptHi}{$-$0.07}
\newcommand{\LKernelSHAPMatFiftyKeptLo}{$-$0.23}
\newcommand{\LKernelSHAPMatFiftyKeptRaw}{0.23}
\newcommand{\LKernelSHAPMatFiftyKeptRawHi}{0.27}
\newcommand{\LKernelSHAPMatFiftyKeptRawLo}{0.19}
\newcommand{\LKernelSHAPMatFiftyLeak}{20.8}
\newcommand{\LKernelSHAPMatFiftyLeakHi}{23.4}
\newcommand{\LKernelSHAPMatFiftyLeakLo}{18.3}
\newcommand{\LKernelSHAPMatFiftyRank}{3.6}
\newcommand{\LKernelSHAPMatFiftyRankHi}{4.3}
\newcommand{\LKernelSHAPMatFiftyRankLo}{2.9}
\newcommand{\LKernelSHAPMatFiftyReloc}{0.19}
\newcommand{\LKernelSHAPMatFiftyRelocHi}{0.26}
\newcommand{\LKernelSHAPMatFiftyRelocLo}{0.12}
\newcommand{\LKernelSHAPMatFiftyRelocN}{5}
\newcommand{\LKernelSHAPMatFiftyRelocTo}{grade}
\newcommand{\LKernelSHAPMatFiftyRelocToN}{4}
\newcommand{\LKernelSHAPMatFiftySigned}{$-$0.08}
\newcommand{\LKernelSHAPMatFiftySignedHi}{$-$0.07}
\newcommand{\LKernelSHAPMatFiftySignedLo}{$-$0.10}
\newcommand{\LKernelSHAPMatFiftyTop}{55.2}
\newcommand{\LKernelSHAPMatFiftyTopHi}{70.0}
\newcommand{\LKernelSHAPMatFiftyTopLo}{40.4}
\newcommand{\LKernelSHAPMatFiftyTopPen}{63.3}
\newcommand{\LKernelSHAPMatFiftyTopPenHi}{74.5}
\newcommand{\LKernelSHAPMatFiftyTopPenLo}{52.1}
\newcommand{\LKernelSHAPMatFiftyTopUnpen}{53.5}
\newcommand{\LKernelSHAPMatFiftyTopUnpenHi}{70.0}
\newcommand{\LKernelSHAPMatFiftyTopUnpenLo}{37.0}
\newcommand{\LKernelSHAPMatGapTotal}{0.85}
\newcommand{\LKernelSHAPMatGapTotalHi}{0.98}
\newcommand{\LKernelSHAPMatGapTotalLo}{0.71}
\newcommand{\LKernelSHAPMatKept}{0.48}
\newcommand{\LKernelSHAPMatKeptHi}{0.63}
\newcommand{\LKernelSHAPMatKeptLo}{0.33}
\newcommand{\LKernelSHAPMatKeptRaw}{0.65}
\newcommand{\LKernelSHAPMatKeptRawHi}{0.74}
\newcommand{\LKernelSHAPMatKeptRawLo}{0.56}
\newcommand{\LKernelSHAPMatLeak}{59.3}
\newcommand{\LKernelSHAPMatLeakHi}{62.4}
\newcommand{\LKernelSHAPMatLeakLo}{56.2}
\newcommand{\LKernelSHAPMatRank}{1.0}
\newcommand{\LKernelSHAPMatRankHi}{1.0}
\newcommand{\LKernelSHAPMatRankLo}{1.0}
\newcommand{\LKernelSHAPMatReloc}{0.19}
\newcommand{\LKernelSHAPMatRelocHi}{0.24}
\newcommand{\LKernelSHAPMatRelocLo}{0.13}
\newcommand{\LKernelSHAPMatRelocN}{5}
\newcommand{\LKernelSHAPMatRelocTo}{age}
\newcommand{\LKernelSHAPMatRelocToN}{2}
\newcommand{\LKernelSHAPMatSeventyFiveAuc}{0.74}
\newcommand{\LKernelSHAPMatSeventyFiveAucHi}{0.75}
\newcommand{\LKernelSHAPMatSeventyFiveAucLo}{0.72}
\newcommand{\LKernelSHAPMatSeventyFiveFid}{74.4}
\newcommand{\LKernelSHAPMatSeventyFiveFidHi}{82.8}
\newcommand{\LKernelSHAPMatSeventyFiveFidLo}{66.0}
\newcommand{\LKernelSHAPMatSeventyFiveFidPen}{67.9}
\newcommand{\LKernelSHAPMatSeventyFiveFidPenHi}{80.5}
\newcommand{\LKernelSHAPMatSeventyFiveFidPenLo}{55.2}
\newcommand{\LKernelSHAPMatSeventyFiveGapTotal}{0.69}
\newcommand{\LKernelSHAPMatSeventyFiveGapTotalHi}{0.80}
\newcommand{\LKernelSHAPMatSeventyFiveGapTotalLo}{0.58}
\newcommand{\LKernelSHAPMatSeventyFiveKept}{0.17}
\newcommand{\LKernelSHAPMatSeventyFiveKeptHi}{0.31}
\newcommand{\LKernelSHAPMatSeventyFiveKeptLo}{0.03}
\newcommand{\LKernelSHAPMatSeventyFiveKeptRaw}{0.44}
\newcommand{\LKernelSHAPMatSeventyFiveKeptRawHi}{0.52}
\newcommand{\LKernelSHAPMatSeventyFiveKeptRawLo}{0.37}
\newcommand{\LKernelSHAPMatSeventyFiveLeak}{38.4}
\newcommand{\LKernelSHAPMatSeventyFiveLeakHi}{42.6}
\newcommand{\LKernelSHAPMatSeventyFiveLeakLo}{34.3}
\newcommand{\LKernelSHAPMatSeventyFiveRank}{1.0}
\newcommand{\LKernelSHAPMatSeventyFiveRankHi}{1.0}
\newcommand{\LKernelSHAPMatSeventyFiveRankLo}{1.0}
\newcommand{\LKernelSHAPMatSeventyFiveReloc}{0.23}
\newcommand{\LKernelSHAPMatSeventyFiveRelocHi}{0.28}
\newcommand{\LKernelSHAPMatSeventyFiveRelocLo}{0.18}
\newcommand{\LKernelSHAPMatSeventyFiveRelocN}{5}
\newcommand{\LKernelSHAPMatSeventyFiveRelocTo}{grade}
\newcommand{\LKernelSHAPMatSeventyFiveRelocToN}{3}
\newcommand{\LKernelSHAPMatSeventyFiveSigned}{$-$0.16}
\newcommand{\LKernelSHAPMatSeventyFiveSignedHi}{$-$0.14}
\newcommand{\LKernelSHAPMatSeventyFiveSignedLo}{$-$0.17}
\newcommand{\LKernelSHAPMatSeventyFiveTop}{82.1}
\newcommand{\LKernelSHAPMatSeventyFiveTopHi}{92.6}
\newcommand{\LKernelSHAPMatSeventyFiveTopLo}{71.6}
\newcommand{\LKernelSHAPMatSeventyFiveTopPen}{89.8}
\newcommand{\LKernelSHAPMatSeventyFiveTopPenHi}{96.6}
\newcommand{\LKernelSHAPMatSeventyFiveTopPenLo}{83.0}
\newcommand{\LKernelSHAPMatSeventyFiveTopUnpen}{80.7}
\newcommand{\LKernelSHAPMatSeventyFiveTopUnpenHi}{92.5}
\newcommand{\LKernelSHAPMatSeventyFiveTopUnpenLo}{68.9}
\newcommand{\LKernelSHAPMatSigned}{$-$0.23}
\newcommand{\LKernelSHAPMatSignedHi}{$-$0.21}
\newcommand{\LKernelSHAPMatSignedLo}{$-$0.26}
\newcommand{\LKernelSHAPMatTop}{95.8}
\newcommand{\LKernelSHAPMatTopHi}{99.5}
\newcommand{\LKernelSHAPMatTopLo}{92.1}
\newcommand{\LKernelSHAPMatTopPen}{98.1}
\newcommand{\LKernelSHAPMatTopPenHi}{103.4}
\newcommand{\LKernelSHAPMatTopPenLo}{92.7}
\newcommand{\LKernelSHAPMatTopUnpen}{95.4}
\newcommand{\LKernelSHAPMatTopUnpenHi}{99.1}
\newcommand{\LKernelSHAPMatTopUnpenLo}{91.7}
\newcommand{\LLIMEBiasedKept}{1.00}
\newcommand{\LLIMEBiasedKeptHi}{1.00}
\newcommand{\LLIMEBiasedKeptLo}{1.00}
\newcommand{\LLIMEBiasedKeptRaw}{1.00}
\newcommand{\LLIMEBiasedKeptRawHi}{1.00}
\newcommand{\LLIMEBiasedKeptRawLo}{1.00}
\newcommand{\LLIMEBiasedRank}{1.0}
\newcommand{\LLIMEBiasedRankHi}{1.0}
\newcommand{\LLIMEBiasedRankLo}{1.0}
\newcommand{\LLIMEBiasedSigned}{$-$0.17}
\newcommand{\LLIMEBiasedSignedHi}{$-$0.12}
\newcommand{\LLIMEBiasedSignedLo}{$-$0.22}
\newcommand{\LLIMEBiasedTop}{94.6}
\newcommand{\LLIMEBiasedTopHi}{99.6}
\newcommand{\LLIMEBiasedTopLo}{89.6}
\newcommand{\LLIMEBiasedTopPen}{93.9}
\newcommand{\LLIMEBiasedTopPenHi}{100.5}
\newcommand{\LLIMEBiasedTopPenLo}{87.3}
\newcommand{\LLIMEBiasedTopUnpen}{94.8}
\newcommand{\LLIMEBiasedTopUnpenHi}{99.5}
\newcommand{\LLIMEBiasedTopUnpenLo}{90.0}
\newcommand{\LLIMEGenAuc}{0.97}
\newcommand{\LLIMEGenAucHi}{0.98}
\newcommand{\LLIMEGenAucLo}{0.96}
\newcommand{\LLIMEGenFid}{100.0}
\newcommand{\LLIMEGenFidHi}{100.0}
\newcommand{\LLIMEGenFidLo}{100.0}
\newcommand{\LLIMEGenFidPen}{100.0}
\newcommand{\LLIMEGenFidPenHi}{100.0}
\newcommand{\LLIMEGenFidPenLo}{100.0}
\newcommand{\LLIMEGenGapTotal}{0.97}
\newcommand{\LLIMEGenGapTotalHi}{1.12}
\newcommand{\LLIMEGenGapTotalLo}{0.81}
\newcommand{\LLIMEGenKept}{0.90}
\newcommand{\LLIMEGenKeptHi}{0.97}
\newcommand{\LLIMEGenKeptLo}{0.83}
\newcommand{\LLIMEGenKeptRaw}{0.96}
\newcommand{\LLIMEGenKeptRawHi}{0.99}
\newcommand{\LLIMEGenKeptRawLo}{0.93}
\newcommand{\LLIMEGenLeak}{94.9}
\newcommand{\LLIMEGenLeakHi}{98.1}
\newcommand{\LLIMEGenLeakLo}{91.8}
\newcommand{\LLIMEGenRank}{1.0}
\newcommand{\LLIMEGenRankHi}{1.0}
\newcommand{\LLIMEGenRankLo}{1.0}
\newcommand{\LLIMEGenReloc}{0.01}
\newcommand{\LLIMEGenRelocHi}{0.16}
\newcommand{\LLIMEGenRelocLo}{$-$0.13}
\newcommand{\LLIMEGenRelocN}{5}
\newcommand{\LLIMEGenRelocTo}{view}
\newcommand{\LLIMEGenRelocToN}{2}
\newcommand{\LLIMEGenSigned}{$-$0.16}
\newcommand{\LLIMEGenSignedHi}{$-$0.12}
\newcommand{\LLIMEGenSignedLo}{$-$0.21}
\newcommand{\LLIMEGenTop}{90.3}
\newcommand{\LLIMEGenTopHi}{99.5}
\newcommand{\LLIMEGenTopLo}{81.1}
\newcommand{\LLIMEGenTopPen}{88.9}
\newcommand{\LLIMEGenTopPenHi}{100.9}
\newcommand{\LLIMEGenTopPenLo}{77.0}
\newcommand{\LLIMEGenTopUnpen}{90.5}
\newcommand{\LLIMEGenTopUnpenHi}{99.2}
\newcommand{\LLIMEGenTopUnpenLo}{81.9}
\newcommand{\LLIMEHonestGapTotal}{0.45}
\newcommand{\LLIMEHonestGapTotalHi}{0.61}
\newcommand{\LLIMEHonestGapTotalLo}{0.29}
\newcommand{\LLIMEHonestKept}{0.00}
\newcommand{\LLIMEHonestKeptHi}{0.00}
\newcommand{\LLIMEHonestKeptLo}{0.00}
\newcommand{\LLIMEHonestKeptRaw}{0.57}
\newcommand{\LLIMEHonestKeptRawHi}{0.65}
\newcommand{\LLIMEHonestKeptRawLo}{0.49}
\newcommand{\LLIMEHonestRank}{1.0}
\newcommand{\LLIMEHonestRankHi}{1.0}
\newcommand{\LLIMEHonestRankLo}{1.0}
\newcommand{\LLIMEHonestReloc}{$-$0.11}
\newcommand{\LLIMEHonestRelocHi}{0.01}
\newcommand{\LLIMEHonestRelocLo}{$-$0.23}
\newcommand{\LLIMEHonestRelocN}{5}
\newcommand{\LLIMEHonestRelocTo}{condition}
\newcommand{\LLIMEHonestRelocToN}{2}
\newcommand{\LLIMEHonestSigned}{$-$0.10}
\newcommand{\LLIMEHonestSignedHi}{$-$0.06}
\newcommand{\LLIMEHonestSignedLo}{$-$0.14}
\newcommand{\LLIMEHonestTop}{66.1}
\newcommand{\LLIMEHonestTopHi}{77.0}
\newcommand{\LLIMEHonestTopLo}{55.2}
\newcommand{\LLIMEHonestTopPen}{68.6}
\newcommand{\LLIMEHonestTopPenHi}{84.4}
\newcommand{\LLIMEHonestTopPenLo}{52.8}
\newcommand{\LLIMEHonestTopUnpen}{65.8}
\newcommand{\LLIMEHonestTopUnpenHi}{76.5}
\newcommand{\LLIMEHonestTopUnpenLo}{55.1}
\newcommand{\LLIMEMatAuc}{1.00}
\newcommand{\LLIMEMatAucHi}{1.00}
\newcommand{\LLIMEMatAucLo}{1.00}
\newcommand{\LLIMEMatFid}{96.9}
\newcommand{\LLIMEMatFidHi}{98.4}
\newcommand{\LLIMEMatFidLo}{95.4}
\newcommand{\LLIMEMatFidPen}{98.9}
\newcommand{\LLIMEMatFidPenHi}{101.9}
\newcommand{\LLIMEMatFidPenLo}{96.0}
\newcommand{\LLIMEMatGapTotal}{0.44}
\newcommand{\LLIMEMatGapTotalHi}{0.78}
\newcommand{\LLIMEMatGapTotalLo}{0.11}
\newcommand{\LLIMEMatKept}{$-$1.24}
\newcommand{\LLIMEMatKeptHi}{$-$0.86}
\newcommand{\LLIMEMatKeptLo}{$-$1.62}
\newcommand{\LLIMEMatKeptRaw}{0.05}
\newcommand{\LLIMEMatKeptRawHi}{0.06}
\newcommand{\LLIMEMatKeptRawLo}{0.04}
\newcommand{\LLIMEMatLeak}{0.0}
\newcommand{\LLIMEMatLeakHi}{0.0}
\newcommand{\LLIMEMatLeakLo}{0.0}
\newcommand{\LLIMEMatRank}{8.8}
\newcommand{\LLIMEMatRankHi}{9.4}
\newcommand{\LLIMEMatRankLo}{8.2}
\newcommand{\LLIMEMatReloc}{0.40}
\newcommand{\LLIMEMatRelocHi}{0.73}
\newcommand{\LLIMEMatRelocLo}{0.06}
\newcommand{\LLIMEMatRelocN}{5}
\newcommand{\LLIMEMatRelocTo}{grade}
\newcommand{\LLIMEMatRelocToN}{4}
\newcommand{\LLIMEMatSigned}{$-$0.01}
\newcommand{\LLIMEMatSignedHi}{$-$0.01}
\newcommand{\LLIMEMatSignedLo}{$-$0.01}
\newcommand{\LLIMEMatTop}{0.8}
\newcommand{\LLIMEMatTopHi}{2.1}
\newcommand{\LLIMEMatTopLo}{$-$0.5}
\newcommand{\LLIMEMatTopPen}{1.6}
\newcommand{\LLIMEMatTopPenHi}{6.0}
\newcommand{\LLIMEMatTopPenLo}{$-$2.8}
\newcommand{\LLIMEMatTopUnpen}{0.6}
\newcommand{\LLIMEMatTopUnpenHi}{1.3}
\newcommand{\LLIMEMatTopUnpenLo}{$-$0.1}
\newcommand{\MechFeatureGeoShapleyFidFifty}{52.1}
\newcommand{\MechFeatureGeoShapleyFidFiftyHi}{55.4}
\newcommand{\MechFeatureGeoShapleyFidFiftyLo}{48.8}
\newcommand{\MechFeatureGeoShapleyFidNinety}{89.7}
\newcommand{\MechFeatureGeoShapleyFidNinetyHi}{91.8}
\newcommand{\MechFeatureGeoShapleyFidNinetyLo}{87.6}
\newcommand{\MechFeatureGeoShapleyFidSeventyFive}{75.0}
\newcommand{\MechFeatureGeoShapleyFidSeventyFiveHi}{80.8}
\newcommand{\MechFeatureGeoShapleyFidSeventyFiveLo}{69.2}
\newcommand{\MechFeatureGeoShapleyM}{10}
\newcommand{\MechFeatureGeoShapleyMAE}{0.003}
\newcommand{\MechFeatureGeoShapleyMeasFifty}{0.29}
\newcommand{\MechFeatureGeoShapleyMeasFiftyHi}{0.32}
\newcommand{\MechFeatureGeoShapleyMeasFiftyLo}{0.26}
\newcommand{\MechFeatureGeoShapleyMeasNinety}{0.72}
\newcommand{\MechFeatureGeoShapleyMeasNinetyHi}{0.79}
\newcommand{\MechFeatureGeoShapleyMeasNinetyLo}{0.65}
\newcommand{\MechFeatureGeoShapleyMeasSeventyFive}{0.50}
\newcommand{\MechFeatureGeoShapleyMeasSeventyFiveHi}{0.53}
\newcommand{\MechFeatureGeoShapleyMeasSeventyFiveLo}{0.47}
\newcommand{\MechFeatureGeoShapleyPredFifty}{0.29}
\newcommand{\MechFeatureGeoShapleyPredFiftyHi}{0.31}
\newcommand{\MechFeatureGeoShapleyPredFiftyLo}{0.26}
\newcommand{\MechFeatureGeoShapleyPredNinety}{0.72}
\newcommand{\MechFeatureGeoShapleyPredNinetyHi}{0.79}
\newcommand{\MechFeatureGeoShapleyPredNinetyLo}{0.65}
\newcommand{\MechFeatureGeoShapleyPredSeventyFive}{0.50}
\newcommand{\MechFeatureGeoShapleyPredSeventyFiveHi}{0.53}
\newcommand{\MechFeatureGeoShapleyPredSeventyFiveLo}{0.46}
\newcommand{\MechFeatureGeoShapleyRhoFifty}{0.28}
\newcommand{\MechFeatureGeoShapleyRhoFiftyHi}{0.30}
\newcommand{\MechFeatureGeoShapleyRhoFiftyLo}{0.26}
\newcommand{\MechFeatureGeoShapleyRhoFirst}{0.91}
\newcommand{\MechFeatureGeoShapleyRhoLast}{0.88}
\newcommand{\MechFeatureGeoShapleyRhoMAE}{0.026}
\newcommand{\MechFeatureGeoShapleyRhoMin}{0.53}
\newcommand{\MechFeatureGeoShapleyRhoMinK}{5}
\newcommand{\MechFeatureGeoShapleyRhoNinety}{0.68}
\newcommand{\MechFeatureGeoShapleyRhoNinetyHi}{0.74}
\newcommand{\MechFeatureGeoShapleyRhoNinetyLo}{0.62}
\newcommand{\MechFeatureGeoShapleyRhoSeventyFive}{0.47}
\newcommand{\MechFeatureGeoShapleyRhoSeventyFiveHi}{0.51}
\newcommand{\MechFeatureGeoShapleyRhoSeventyFiveLo}{0.43}
\newcommand{\MechFeatureGeoShapleySD}{0.035}
\newcommand{\MechFeatureKernelSHAPFidFifty}{50.5}
\newcommand{\MechFeatureKernelSHAPFidFiftyHi}{55.3}
\newcommand{\MechFeatureKernelSHAPFidFiftyLo}{45.7}
\newcommand{\MechFeatureKernelSHAPFidNinety}{89.5}
\newcommand{\MechFeatureKernelSHAPFidNinetyHi}{91.5}
\newcommand{\MechFeatureKernelSHAPFidNinetyLo}{87.5}
\newcommand{\MechFeatureKernelSHAPFidSeventyFive}{74.9}
\newcommand{\MechFeatureKernelSHAPFidSeventyFiveHi}{79.7}
\newcommand{\MechFeatureKernelSHAPFidSeventyFiveLo}{70.1}
\newcommand{\MechFeatureKernelSHAPM}{11}
\newcommand{\MechFeatureKernelSHAPMAE}{0.002}
\newcommand{\MechFeatureKernelSHAPMeasFifty}{0.24}
\newcommand{\MechFeatureKernelSHAPMeasFiftyHi}{0.26}
\newcommand{\MechFeatureKernelSHAPMeasFiftyLo}{0.22}
\newcommand{\MechFeatureKernelSHAPMeasNinety}{0.67}
\newcommand{\MechFeatureKernelSHAPMeasNinetyHi}{0.74}
\newcommand{\MechFeatureKernelSHAPMeasNinetyLo}{0.60}
\newcommand{\MechFeatureKernelSHAPMeasSeventyFive}{0.47}
\newcommand{\MechFeatureKernelSHAPMeasSeventyFiveHi}{0.52}
\newcommand{\MechFeatureKernelSHAPMeasSeventyFiveLo}{0.42}
\newcommand{\MechFeatureKernelSHAPPredFifty}{0.24}
\newcommand{\MechFeatureKernelSHAPPredFiftyHi}{0.26}
\newcommand{\MechFeatureKernelSHAPPredFiftyLo}{0.22}
\newcommand{\MechFeatureKernelSHAPPredNinety}{0.67}
\newcommand{\MechFeatureKernelSHAPPredNinetyHi}{0.74}
\newcommand{\MechFeatureKernelSHAPPredNinetyLo}{0.60}
\newcommand{\MechFeatureKernelSHAPPredSeventyFive}{0.47}
\newcommand{\MechFeatureKernelSHAPPredSeventyFiveHi}{0.52}
\newcommand{\MechFeatureKernelSHAPPredSeventyFiveLo}{0.42}
\newcommand{\MechFeatureKernelSHAPRhoFifty}{0.25}
\newcommand{\MechFeatureKernelSHAPRhoFiftyHi}{0.27}
\newcommand{\MechFeatureKernelSHAPRhoFiftyLo}{0.22}
\newcommand{\MechFeatureKernelSHAPRhoFirst}{0.87}
\newcommand{\MechFeatureKernelSHAPRhoLast}{0.86}
\newcommand{\MechFeatureKernelSHAPRhoMAE}{0.019}
\newcommand{\MechFeatureKernelSHAPRhoMin}{0.48}
\newcommand{\MechFeatureKernelSHAPRhoMinK}{6}
\newcommand{\MechFeatureKernelSHAPRhoNinety}{0.64}
\newcommand{\MechFeatureKernelSHAPRhoNinetyHi}{0.71}
\newcommand{\MechFeatureKernelSHAPRhoNinetyLo}{0.57}
\newcommand{\MechFeatureKernelSHAPRhoSeventyFive}{0.45}
\newcommand{\MechFeatureKernelSHAPRhoSeventyFiveHi}{0.50}
\newcommand{\MechFeatureKernelSHAPRhoSeventyFiveLo}{0.40}
\newcommand{\MechFeatureKernelSHAPSD}{0.037}
\newcommand{\MechLocationGeoShapleyFidFifty}{52.1}
\newcommand{\MechLocationGeoShapleyFidFiftyHi}{57.6}
\newcommand{\MechLocationGeoShapleyFidFiftyLo}{46.6}
\newcommand{\MechLocationGeoShapleyFidNinety}{88.8}
\newcommand{\MechLocationGeoShapleyFidNinetyHi}{92.8}
\newcommand{\MechLocationGeoShapleyFidNinetyLo}{84.8}
\newcommand{\MechLocationGeoShapleyFidSeventyFive}{74.5}
\newcommand{\MechLocationGeoShapleyFidSeventyFiveHi}{78.9}
\newcommand{\MechLocationGeoShapleyFidSeventyFiveLo}{70.1}
\newcommand{\MechLocationGeoShapleyM}{9}
\newcommand{\MechLocationGeoShapleyMAE}{0.002}
\newcommand{\MechLocationGeoShapleyMeasFifty}{0.21}
\newcommand{\MechLocationGeoShapleyMeasFiftyHi}{0.26}
\newcommand{\MechLocationGeoShapleyMeasFiftyLo}{0.16}
\newcommand{\MechLocationGeoShapleyMeasNinety}{0.65}
\newcommand{\MechLocationGeoShapleyMeasNinetyHi}{0.72}
\newcommand{\MechLocationGeoShapleyMeasNinetyLo}{0.57}
\newcommand{\MechLocationGeoShapleyMeasSeventyFive}{0.42}
\newcommand{\MechLocationGeoShapleyMeasSeventyFiveHi}{0.50}
\newcommand{\MechLocationGeoShapleyMeasSeventyFiveLo}{0.33}
\newcommand{\MechLocationGeoShapleyPredFifty}{0.21}
\newcommand{\MechLocationGeoShapleyPredFiftyHi}{0.27}
\newcommand{\MechLocationGeoShapleyPredFiftyLo}{0.16}
\newcommand{\MechLocationGeoShapleyPredNinety}{0.65}
\newcommand{\MechLocationGeoShapleyPredNinetyHi}{0.72}
\newcommand{\MechLocationGeoShapleyPredNinetyLo}{0.57}
\newcommand{\MechLocationGeoShapleyPredSeventyFive}{0.42}
\newcommand{\MechLocationGeoShapleyPredSeventyFiveHi}{0.50}
\newcommand{\MechLocationGeoShapleyPredSeventyFiveLo}{0.33}
\newcommand{\MechLocationGeoShapleyRhoFifty}{0.26}
\newcommand{\MechLocationGeoShapleyRhoFiftyHi}{0.31}
\newcommand{\MechLocationGeoShapleyRhoFiftyLo}{0.21}
\newcommand{\MechLocationGeoShapleyRhoFirst}{0.84}
\newcommand{\MechLocationGeoShapleyRhoLast}{0.82}
\newcommand{\MechLocationGeoShapleyRhoMAE}{0.040}
\newcommand{\MechLocationGeoShapleyRhoMin}{0.58}
\newcommand{\MechLocationGeoShapleyRhoMinK}{5}
\newcommand{\MechLocationGeoShapleyRhoNinety}{0.68}
\newcommand{\MechLocationGeoShapleyRhoNinetyHi}{0.75}
\newcommand{\MechLocationGeoShapleyRhoNinetyLo}{0.61}
\newcommand{\MechLocationGeoShapleyRhoSeventyFive}{0.46}
\newcommand{\MechLocationGeoShapleyRhoSeventyFiveHi}{0.55}
\newcommand{\MechLocationGeoShapleyRhoSeventyFiveLo}{0.38}
\newcommand{\MechLocationGeoShapleySD}{0.057}
\newcommand{\MinutesPerSeed}{36}
\newcommand{\NKernelSHAPBiasedKept}{1.00}
\newcommand{\NKernelSHAPBiasedKeptHi}{1.00}
\newcommand{\NKernelSHAPBiasedKeptLo}{1.00}
\newcommand{\NKernelSHAPBiasedRank}{1.0}
\newcommand{\NKernelSHAPBiasedRankHi}{1.0}
\newcommand{\NKernelSHAPBiasedRankLo}{1.0}
\newcommand{\NKernelSHAPBiasedSigned}{$-$0.24}
\newcommand{\NKernelSHAPBiasedSignedHi}{$-$0.23}
\newcommand{\NKernelSHAPBiasedSignedLo}{$-$0.26}
\newcommand{\NKernelSHAPBiasedTop}{80.2}
\newcommand{\NKernelSHAPBiasedTopHi}{93.0}
\newcommand{\NKernelSHAPBiasedTopLo}{67.4}
\newcommand{\NKernelSHAPBiasedTopPen}{100.0}
\newcommand{\NKernelSHAPBiasedTopPenHi}{100.0}
\newcommand{\NKernelSHAPBiasedTopPenLo}{100.0}
\newcommand{\NKernelSHAPBiasedTopUnpen}{76.3}
\newcommand{\NKernelSHAPBiasedTopUnpenHi}{91.5}
\newcommand{\NKernelSHAPBiasedTopUnpenLo}{61.1}
\newcommand{\NKernelSHAPGenAuc}{0.68}
\newcommand{\NKernelSHAPGenAucHi}{0.70}
\newcommand{\NKernelSHAPGenAucLo}{0.66}
\newcommand{\NKernelSHAPGenFid}{76.5}
\newcommand{\NKernelSHAPGenFidHi}{80.2}
\newcommand{\NKernelSHAPGenFidLo}{72.8}
\newcommand{\NKernelSHAPGenFidPen}{78.1}
\newcommand{\NKernelSHAPGenFidPenHi}{84.0}
\newcommand{\NKernelSHAPGenFidPenLo}{72.2}
\newcommand{\NKernelSHAPGenGapTotal}{0.79}
\newcommand{\NKernelSHAPGenGapTotalHi}{0.86}
\newcommand{\NKernelSHAPGenGapTotalLo}{0.72}
\newcommand{\NKernelSHAPGenKept}{0.61}
\newcommand{\NKernelSHAPGenKeptHi}{0.63}
\newcommand{\NKernelSHAPGenKeptLo}{0.59}
\newcommand{\NKernelSHAPGenLeak}{50.1}
\newcommand{\NKernelSHAPGenLeakHi}{52.6}
\newcommand{\NKernelSHAPGenLeakLo}{47.6}
\newcommand{\NKernelSHAPGenRank}{3.2}
\newcommand{\NKernelSHAPGenRankHi}{4.2}
\newcommand{\NKernelSHAPGenRankLo}{2.2}
\newcommand{\NKernelSHAPGenReloc}{0.17}
\newcommand{\NKernelSHAPGenRelocHi}{0.24}
\newcommand{\NKernelSHAPGenRelocLo}{0.09}
\newcommand{\NKernelSHAPGenRelocN}{5}
\newcommand{\NKernelSHAPGenRelocTo}{grade}
\newcommand{\NKernelSHAPGenRelocToN}{4}
\newcommand{\NKernelSHAPGenSigned}{$-$0.15}
\newcommand{\NKernelSHAPGenSignedHi}{$-$0.14}
\newcommand{\NKernelSHAPGenSignedLo}{$-$0.16}
\newcommand{\NKernelSHAPGenTop}{52.4}
\newcommand{\NKernelSHAPGenTopHi}{68.8}
\newcommand{\NKernelSHAPGenTopLo}{36.0}
\newcommand{\NKernelSHAPGenTopPen}{98.2}
\newcommand{\NKernelSHAPGenTopPenHi}{100.3}
\newcommand{\NKernelSHAPGenTopPenLo}{96.0}
\newcommand{\NKernelSHAPGenTopUnpen}{43.4}
\newcommand{\NKernelSHAPGenTopUnpenHi}{63.8}
\newcommand{\NKernelSHAPGenTopUnpenLo}{23.0}
\newcommand{\NKernelSHAPHonestGapTotal}{0.00}
\newcommand{\NKernelSHAPHonestGapTotalHi}{0.00}
\newcommand{\NKernelSHAPHonestGapTotalLo}{0.00}
\newcommand{\NKernelSHAPHonestKept}{0.00}
\newcommand{\NKernelSHAPHonestKeptHi}{0.00}
\newcommand{\NKernelSHAPHonestKeptLo}{0.00}
\newcommand{\NKernelSHAPHonestRank}{9.0}
\newcommand{\NKernelSHAPHonestRankHi}{9.0}
\newcommand{\NKernelSHAPHonestRankLo}{9.0}
\newcommand{\NKernelSHAPHonestReloc}{0.00}
\newcommand{\NKernelSHAPHonestRelocHi}{0.00}
\newcommand{\NKernelSHAPHonestRelocLo}{0.00}
\newcommand{\NKernelSHAPHonestRelocN}{5}
\newcommand{\NKernelSHAPHonestRelocTo}{waterfront}
\newcommand{\NKernelSHAPHonestRelocToN}{2}
\newcommand{\NKernelSHAPHonestSigned}{0.00}
\newcommand{\NKernelSHAPHonestSignedHi}{0.00}
\newcommand{\NKernelSHAPHonestSignedLo}{0.00}
\newcommand{\NKernelSHAPHonestTop}{0.0}
\newcommand{\NKernelSHAPHonestTopHi}{0.0}
\newcommand{\NKernelSHAPHonestTopLo}{0.0}
\newcommand{\NKernelSHAPHonestTopPen}{0.0}
\newcommand{\NKernelSHAPHonestTopPenHi}{0.0}
\newcommand{\NKernelSHAPHonestTopPenLo}{0.0}
\newcommand{\NKernelSHAPHonestTopUnpen}{0.0}
\newcommand{\NKernelSHAPHonestTopUnpenHi}{0.0}
\newcommand{\NKernelSHAPHonestTopUnpenLo}{0.0}
\newcommand{\NKernelSHAPMatAuc}{0.70}
\newcommand{\NKernelSHAPMatAucHi}{0.71}
\newcommand{\NKernelSHAPMatAucLo}{0.69}
\newcommand{\NKernelSHAPMatFid}{90.2}
\newcommand{\NKernelSHAPMatFidHi}{92.2}
\newcommand{\NKernelSHAPMatFidLo}{88.2}
\newcommand{\NKernelSHAPMatFidPen}{88.5}
\newcommand{\NKernelSHAPMatFidPenHi}{94.4}
\newcommand{\NKernelSHAPMatFidPenLo}{82.7}
\newcommand{\NKernelSHAPMatGapTotal}{0.88}
\newcommand{\NKernelSHAPMatGapTotalHi}{0.95}
\newcommand{\NKernelSHAPMatGapTotalLo}{0.82}
\newcommand{\NKernelSHAPMatKept}{0.74}
\newcommand{\NKernelSHAPMatKeptHi}{0.78}
\newcommand{\NKernelSHAPMatKeptLo}{0.70}
\newcommand{\NKernelSHAPMatLeak}{66.0}
\newcommand{\NKernelSHAPMatLeakHi}{68.0}
\newcommand{\NKernelSHAPMatLeakLo}{64.0}
\newcommand{\NKernelSHAPMatRank}{3.0}
\newcommand{\NKernelSHAPMatRankHi}{3.9}
\newcommand{\NKernelSHAPMatRankLo}{2.1}
\newcommand{\NKernelSHAPMatReloc}{0.14}
\newcommand{\NKernelSHAPMatRelocHi}{0.18}
\newcommand{\NKernelSHAPMatRelocLo}{0.09}
\newcommand{\NKernelSHAPMatRelocN}{5}
\newcommand{\NKernelSHAPMatRelocTo}{grade}
\newcommand{\NKernelSHAPMatRelocToN}{3}
\newcommand{\NKernelSHAPMatSigned}{$-$0.18}
\newcommand{\NKernelSHAPMatSignedHi}{$-$0.16}
\newcommand{\NKernelSHAPMatSignedLo}{$-$0.20}
\newcommand{\NKernelSHAPMatTop}{60.6}
\newcommand{\NKernelSHAPMatTopHi}{79.2}
\newcommand{\NKernelSHAPMatTopLo}{42.0}
\newcommand{\NKernelSHAPMatTopPen}{100.0}
\newcommand{\NKernelSHAPMatTopPenHi}{100.0}
\newcommand{\NKernelSHAPMatTopPenLo}{100.0}
\newcommand{\NKernelSHAPMatTopUnpen}{52.8}
\newcommand{\NKernelSHAPMatTopUnpenHi}{75.1}
\newcommand{\NKernelSHAPMatTopUnpenLo}{30.6}
\newcommand{\NLIMEBiasedKept}{1.00}
\newcommand{\NLIMEBiasedKeptHi}{1.00}
\newcommand{\NLIMEBiasedKeptLo}{1.00}
\newcommand{\NLIMEBiasedRank}{2.8}
\newcommand{\NLIMEBiasedRankHi}{3.4}
\newcommand{\NLIMEBiasedRankLo}{2.2}
\newcommand{\NLIMEBiasedSigned}{$-$0.24}
\newcommand{\NLIMEBiasedSignedHi}{$-$0.23}
\newcommand{\NLIMEBiasedSignedLo}{$-$0.24}
\newcommand{\NLIMEBiasedTop}{54.7}
\newcommand{\NLIMEBiasedTopHi}{63.9}
\newcommand{\NLIMEBiasedTopLo}{45.5}
\newcommand{\NLIMEBiasedTopPen}{100.0}
\newcommand{\NLIMEBiasedTopPenHi}{100.0}
\newcommand{\NLIMEBiasedTopPenLo}{100.0}
\newcommand{\NLIMEBiasedTopUnpen}{46.1}
\newcommand{\NLIMEBiasedTopUnpenHi}{55.3}
\newcommand{\NLIMEBiasedTopUnpenLo}{37.0}
\newcommand{\NLIMEGenAuc}{0.94}
\newcommand{\NLIMEGenAucHi}{0.96}
\newcommand{\NLIMEGenAucLo}{0.93}
\newcommand{\NLIMEGenFid}{100.0}
\newcommand{\NLIMEGenFidHi}{100.0}
\newcommand{\NLIMEGenFidLo}{100.0}
\newcommand{\NLIMEGenFidPen}{100.0}
\newcommand{\NLIMEGenFidPenHi}{100.0}
\newcommand{\NLIMEGenFidPenLo}{100.0}
\newcommand{\NLIMEGenGapTotal}{1.00}
\newcommand{\NLIMEGenGapTotalHi}{1.04}
\newcommand{\NLIMEGenGapTotalLo}{0.96}
\newcommand{\NLIMEGenKept}{0.99}
\newcommand{\NLIMEGenKeptHi}{1.01}
\newcommand{\NLIMEGenKeptLo}{0.97}
\newcommand{\NLIMEGenLeak}{98.8}
\newcommand{\NLIMEGenLeakHi}{100.8}
\newcommand{\NLIMEGenLeakLo}{96.7}
\newcommand{\NLIMEGenRank}{2.8}
\newcommand{\NLIMEGenRankHi}{3.4}
\newcommand{\NLIMEGenRankLo}{2.2}
\newcommand{\NLIMEGenReloc}{0.01}
\newcommand{\NLIMEGenRelocHi}{0.03}
\newcommand{\NLIMEGenRelocLo}{$-$0.01}
\newcommand{\NLIMEGenRelocN}{5}
\newcommand{\NLIMEGenRelocTo}{waterfront}
\newcommand{\NLIMEGenRelocToN}{4}
\newcommand{\NLIMEGenSigned}{$-$0.24}
\newcommand{\NLIMEGenSignedHi}{$-$0.23}
\newcommand{\NLIMEGenSignedLo}{$-$0.24}
\newcommand{\NLIMEGenTop}{51.7}
\newcommand{\NLIMEGenTopHi}{61.8}
\newcommand{\NLIMEGenTopLo}{41.6}
\newcommand{\NLIMEGenTopPen}{100.0}
\newcommand{\NLIMEGenTopPenHi}{100.0}
\newcommand{\NLIMEGenTopPenLo}{100.0}
\newcommand{\NLIMEGenTopUnpen}{42.6}
\newcommand{\NLIMEGenTopUnpenHi}{52.8}
\newcommand{\NLIMEGenTopUnpenLo}{32.4}
\newcommand{\NLIMEHonestGapTotal}{$-$0.01}
\newcommand{\NLIMEHonestGapTotalHi}{0.02}
\newcommand{\NLIMEHonestGapTotalLo}{$-$0.04}
\newcommand{\NLIMEHonestKept}{0.00}
\newcommand{\NLIMEHonestKeptHi}{0.01}
\newcommand{\NLIMEHonestKeptLo}{$-$0.01}
\newcommand{\NLIMEHonestRank}{9.0}
\newcommand{\NLIMEHonestRankHi}{9.0}
\newcommand{\NLIMEHonestRankLo}{9.0}
\newcommand{\NLIMEHonestReloc}{$-$0.01}
\newcommand{\NLIMEHonestRelocHi}{0.02}
\newcommand{\NLIMEHonestRelocLo}{$-$0.04}
\newcommand{\NLIMEHonestRelocN}{5}
\newcommand{\NLIMEHonestRelocTo}{waterfront}
\newcommand{\NLIMEHonestRelocToN}{4}
\newcommand{\NLIMEHonestSigned}{0.00}
\newcommand{\NLIMEHonestSignedHi}{0.00}
\newcommand{\NLIMEHonestSignedLo}{0.00}
\newcommand{\NLIMEHonestTop}{0.2}
\newcommand{\NLIMEHonestTopHi}{0.8}
\newcommand{\NLIMEHonestTopLo}{$-$0.4}
\newcommand{\NLIMEHonestTopPen}{0.0}
\newcommand{\NLIMEHonestTopPenHi}{0.0}
\newcommand{\NLIMEHonestTopPenLo}{0.0}
\newcommand{\NLIMEHonestTopUnpen}{0.2}
\newcommand{\NLIMEHonestTopUnpenHi}{0.9}
\newcommand{\NLIMEHonestTopUnpenLo}{$-$0.4}
\newcommand{\NLIMEMatAuc}{1.00}
\newcommand{\NLIMEMatAucHi}{1.00}
\newcommand{\NLIMEMatAucLo}{1.00}
\newcommand{\NLIMEMatFid}{96.6}
\newcommand{\NLIMEMatFidHi}{98.4}
\newcommand{\NLIMEMatFidLo}{94.8}
\newcommand{\NLIMEMatFidPen}{97.9}
\newcommand{\NLIMEMatFidPenHi}{102.1}
\newcommand{\NLIMEMatFidPenLo}{93.7}
\newcommand{\NLIMEMatGapTotal}{0.12}
\newcommand{\NLIMEMatGapTotalHi}{0.18}
\newcommand{\NLIMEMatGapTotalLo}{0.06}
\newcommand{\NLIMEMatKept}{0.02}
\newcommand{\NLIMEMatKeptHi}{0.03}
\newcommand{\NLIMEMatKeptLo}{0.00}
\newcommand{\NLIMEMatLeak}{0.0}
\newcommand{\NLIMEMatLeakHi}{0.0}
\newcommand{\NLIMEMatLeakLo}{0.0}
\newcommand{\NLIMEMatRank}{9.0}
\newcommand{\NLIMEMatRankHi}{9.0}
\newcommand{\NLIMEMatRankLo}{9.0}
\newcommand{\NLIMEMatReloc}{0.10}
\newcommand{\NLIMEMatRelocHi}{0.16}
\newcommand{\NLIMEMatRelocLo}{0.05}
\newcommand{\NLIMEMatRelocN}{5}
\newcommand{\NLIMEMatRelocTo}{grade}
\newcommand{\NLIMEMatRelocToN}{2}
\newcommand{\NLIMEMatSigned}{0.00}
\newcommand{\NLIMEMatSignedHi}{0.00}
\newcommand{\NLIMEMatSignedLo}{$-$0.01}
\newcommand{\NLIMEMatTop}{0.4}
\newcommand{\NLIMEMatTopHi}{1.2}
\newcommand{\NLIMEMatTopLo}{$-$0.4}
\newcommand{\NLIMEMatTopPen}{0.0}
\newcommand{\NLIMEMatTopPenHi}{0.0}
\newcommand{\NLIMEMatTopPenLo}{0.0}
\newcommand{\NLIMEMatTopUnpen}{0.5}
\newcommand{\NLIMEMatTopUnpenHi}{1.5}
\newcommand{\NLIMEMatTopUnpenLo}{$-$0.5}
\newcommand{\NPenFeature}{32}
\newcommand{\NPenLocation}{32}
\newcommand{\PenTenGeoShapleyBiasedTopPen}{93.3}
\newcommand{\PenTenGeoShapleyBiasedTopPenHi}{99.3}
\newcommand{\PenTenGeoShapleyBiasedTopPenLo}{87.3}
\newcommand{\PenTenGeoShapleyFiftyKept}{0.39}
\newcommand{\PenTenGeoShapleyFiftyKeptHi}{0.42}
\newcommand{\PenTenGeoShapleyFiftyKeptLo}{0.36}
\newcommand{\PenTenGeoShapleyKept}{0.85}
\newcommand{\PenTenGeoShapleyKeptHi}{0.94}
\newcommand{\PenTenGeoShapleyKeptLo}{0.77}
\newcommand{\PenTenGeoShapleyTopPen}{74.0}
\newcommand{\PenTenGeoShapleyTopPenHi}{80.8}
\newcommand{\PenTenGeoShapleyTopPenLo}{67.2}
\newcommand{\PenTenKernelSHAPBiasedTopPen}{93.3}
\newcommand{\PenTenKernelSHAPBiasedTopPenHi}{99.3}
\newcommand{\PenTenKernelSHAPBiasedTopPenLo}{87.3}
\newcommand{\PenTenKernelSHAPFiftyKept}{0.32}
\newcommand{\PenTenKernelSHAPFiftyKeptHi}{0.35}
\newcommand{\PenTenKernelSHAPFiftyKeptLo}{0.29}
\newcommand{\PenTenKernelSHAPKept}{0.80}
\newcommand{\PenTenKernelSHAPKeptHi}{0.89}
\newcommand{\PenTenKernelSHAPKeptLo}{0.71}
\newcommand{\PenTenKernelSHAPTopPen}{68.1}
\newcommand{\PenTenKernelSHAPTopPenHi}{74.9}
\newcommand{\PenTenKernelSHAPTopPenLo}{61.3}
\newcommand{\PenTenLIMEBiasedTopPen}{55.8}
\newcommand{\PenTenLIMEBiasedTopPenHi}{72.7}
\newcommand{\PenTenLIMEBiasedTopPenLo}{38.8}
\newcommand{\PenTenLIMEFiftyKept}{--}
\newcommand{\PenTenLIMEFiftyKeptHi}{--}
\newcommand{\PenTenLIMEFiftyKeptLo}{--}
\newcommand{\PenTenLIMEKept}{0.03}
\newcommand{\PenTenLIMEKeptHi}{0.08}
\newcommand{\PenTenLIMEKeptLo}{$-$0.02}
\newcommand{\PenTenLIMETopPen}{0.0}
\newcommand{\PenTenLIMETopPenHi}{0.0}
\newcommand{\PenTenLIMETopPenLo}{0.0}
\newcommand{\PsiRtwoFeature}{0.65}
\newcommand{\PsiRtwoFeatureHi}{0.69}
\newcommand{\PsiRtwoFeatureLo}{0.61}
\newcommand{\PsiRtwoForest}{0.65}
\newcommand{\PsiRtwoForestHi}{0.68}
\newcommand{\PsiRtwoForestLo}{0.62}
\newcommand{\PsiRtwoLocation}{0.65}
\newcommand{\PsiRtwoLocationHi}{0.69}
\newcommand{\PsiRtwoLocationLo}{0.61}
\newcommand{\PsiRtwoNet}{0.61}
\newcommand{\PsiRtwoNetHi}{0.65}
\newcommand{\PsiRtwoNetLo}{0.57}
\newcommand{\RasterAgree}{99.97}
\newcommand{\RasterAgreeHi}{100.0}
\newcommand{\RasterAgreeLo}{99.9}
\newcommand{\RobBackgroundCAArowsGeoShapleyKept}{0.72}
\newcommand{\RobBackgroundCAArowsGeoShapleyKeptHi}{0.79}
\newcommand{\RobBackgroundCAArowsGeoShapleyKeptLo}{0.65}
\newcommand{\RobBackgroundCAArowsGeoShapleyLeak}{60.3}
\newcommand{\RobBackgroundCAArowsGeoShapleyLeakHi}{63.9}
\newcommand{\RobBackgroundCAArowsGeoShapleyLeakLo}{56.6}
\newcommand{\RobBackgroundCAArowsGeoShapleyN}{5}
\newcommand{\RobBackgroundCAArowsGeoShapleyTopPen}{100.0}
\newcommand{\RobBackgroundCAArowsGeoShapleyTopPenHi}{100.0}
\newcommand{\RobBackgroundCAArowsGeoShapleyTopPenLo}{100.0}
\newcommand{\RobBackgroundCAArowsKernelSHAPKept}{0.67}
\newcommand{\RobBackgroundCAArowsKernelSHAPKeptHi}{0.74}
\newcommand{\RobBackgroundCAArowsKernelSHAPKeptLo}{0.60}
\newcommand{\RobBackgroundCAArowsKernelSHAPLeak}{53.1}
\newcommand{\RobBackgroundCAArowsKernelSHAPLeakHi}{56.1}
\newcommand{\RobBackgroundCAArowsKernelSHAPLeakLo}{50.1}
\newcommand{\RobBackgroundCAArowsKernelSHAPN}{5}
\newcommand{\RobBackgroundCAArowsKernelSHAPTopPen}{99.5}
\newcommand{\RobBackgroundCAArowsKernelSHAPTopPenHi}{100.9}
\newcommand{\RobBackgroundCAArowsKernelSHAPTopPenLo}{98.1}
\newcommand{\RobBackgroundCFrowsGeoShapleyKept}{0.72}
\newcommand{\RobBackgroundCFrowsGeoShapleyKeptHi}{0.77}
\newcommand{\RobBackgroundCFrowsGeoShapleyKeptLo}{0.66}
\newcommand{\RobBackgroundCFrowsGeoShapleyLeak}{61.0}
\newcommand{\RobBackgroundCFrowsGeoShapleyLeakHi}{64.6}
\newcommand{\RobBackgroundCFrowsGeoShapleyLeakLo}{57.4}
\newcommand{\RobBackgroundCFrowsGeoShapleyN}{5}
\newcommand{\RobBackgroundCFrowsGeoShapleyTopPen}{98.8}
\newcommand{\RobBackgroundCFrowsGeoShapleyTopPenHi}{100.8}
\newcommand{\RobBackgroundCFrowsGeoShapleyTopPenLo}{96.9}
\newcommand{\RobBackgroundCFrowsKernelSHAPKept}{0.67}
\newcommand{\RobBackgroundCFrowsKernelSHAPKeptHi}{0.74}
\newcommand{\RobBackgroundCFrowsKernelSHAPKeptLo}{0.61}
\newcommand{\RobBackgroundCFrowsKernelSHAPLeak}{54.2}
\newcommand{\RobBackgroundCFrowsKernelSHAPLeakHi}{57.6}
\newcommand{\RobBackgroundCFrowsKernelSHAPLeakLo}{50.8}
\newcommand{\RobBackgroundCFrowsKernelSHAPN}{5}
\newcommand{\RobBackgroundCFrowsKernelSHAPTopPen}{98.3}
\newcommand{\RobBackgroundCFrowsKernelSHAPTopPenHi}{101.3}
\newcommand{\RobBackgroundCFrowsKernelSHAPTopPenLo}{95.4}
\newcommand{\RobBreakList}{LIME Gaussian sampler (LIME: partly fooled)}
\newcommand{\RobCells}{14}
\newcommand{\RobConfigs}{7}
\newcommand{\RobHold}{13}
\newcommand{\RobLIMEFAAAsamplesLIMEKept}{0.00}
\newcommand{\RobLIMEFAAAsamplesLIMEKeptHi}{0.01}
\newcommand{\RobLIMEFAAAsamplesLIMEKeptLo}{$-$0.01}
\newcommand{\RobLIMEFAAAsamplesLIMELeak}{0.0}
\newcommand{\RobLIMEFAAAsamplesLIMELeakHi}{0.0}
\newcommand{\RobLIMEFAAAsamplesLIMELeakLo}{0.0}
\newcommand{\RobLIMEFAAAsamplesLIMEN}{5}
\newcommand{\RobLIMEFAAAsamplesLIMETopPen}{0.0}
\newcommand{\RobLIMEFAAAsamplesLIMETopPenHi}{0.0}
\newcommand{\RobLIMEFAAAsamplesLIMETopPenLo}{0.0}
\newcommand{\RobLIMEGaussiansampleLIMEKept}{0.37}
\newcommand{\RobLIMEGaussiansampleLIMEKeptHi}{0.40}
\newcommand{\RobLIMEGaussiansampleLIMEKeptLo}{0.35}
\newcommand{\RobLIMEGaussiansampleLIMELeak}{31.0}
\newcommand{\RobLIMEGaussiansampleLIMELeakHi}{31.8}
\newcommand{\RobLIMEGaussiansampleLIMELeakLo}{30.2}
\newcommand{\RobLIMEGaussiansampleLIMEN}{5}
\newcommand{\RobLIMEGaussiansampleLIMETopPen}{50.1}
\newcommand{\RobLIMEGaussiansampleLIMETopPenHi}{68.8}
\newcommand{\RobLIMEGaussiansampleLIMETopPenLo}{31.4}
\newcommand{\RobMainXGBoostdefaultGeoShapleyKept}{0.72}
\newcommand{\RobMainXGBoostdefaultGeoShapleyKeptHi}{0.79}
\newcommand{\RobMainXGBoostdefaultGeoShapleyKeptLo}{0.65}
\newcommand{\RobMainXGBoostdefaultGeoShapleyLeak}{60.8}
\newcommand{\RobMainXGBoostdefaultGeoShapleyLeakHi}{64.1}
\newcommand{\RobMainXGBoostdefaultGeoShapleyLeakLo}{57.6}
\newcommand{\RobMainXGBoostdefaultGeoShapleyN}{5}
\newcommand{\RobMainXGBoostdefaultGeoShapleyTopPen}{100.0}
\newcommand{\RobMainXGBoostdefaultGeoShapleyTopPenHi}{100.0}
\newcommand{\RobMainXGBoostdefaultGeoShapleyTopPenLo}{100.0}
\newcommand{\RobMainXGBoostdefaultKernelSHAPKept}{0.67}
\newcommand{\RobMainXGBoostdefaultKernelSHAPKeptHi}{0.74}
\newcommand{\RobMainXGBoostdefaultKernelSHAPKeptLo}{0.60}
\newcommand{\RobMainXGBoostdefaultKernelSHAPLeak}{53.6}
\newcommand{\RobMainXGBoostdefaultKernelSHAPLeakHi}{56.7}
\newcommand{\RobMainXGBoostdefaultKernelSHAPLeakLo}{50.6}
\newcommand{\RobMainXGBoostdefaultKernelSHAPN}{5}
\newcommand{\RobMainXGBoostdefaultKernelSHAPTopPen}{99.5}
\newcommand{\RobMainXGBoostdefaultKernelSHAPTopPenHi}{100.9}
\newcommand{\RobMainXGBoostdefaultKernelSHAPTopPenLo}{98.1}
\newcommand{\RobMainXGBoostdefaultLIMEKept}{0.02}
\newcommand{\RobMainXGBoostdefaultLIMEKeptHi}{0.04}
\newcommand{\RobMainXGBoostdefaultLIMEKeptLo}{0.00}
\newcommand{\RobMainXGBoostdefaultLIMELeak}{0.0}
\newcommand{\RobMainXGBoostdefaultLIMELeakHi}{0.0}
\newcommand{\RobMainXGBoostdefaultLIMELeakLo}{0.0}
\newcommand{\RobMainXGBoostdefaultLIMEN}{5}
\newcommand{\RobMainXGBoostdefaultLIMETopPen}{0.0}
\newcommand{\RobMainXGBoostdefaultLIMETopPenHi}{0.0}
\newcommand{\RobMainXGBoostdefaultLIMETopPenLo}{0.0}
\newcommand{\RobNeuralnetGeoShapleyKept}{0.72}
\newcommand{\RobNeuralnetGeoShapleyKeptHi}{0.79}
\newcommand{\RobNeuralnetGeoShapleyKeptLo}{0.65}
\newcommand{\RobNeuralnetGeoShapleyLeak}{60.8}
\newcommand{\RobNeuralnetGeoShapleyLeakHi}{64.1}
\newcommand{\RobNeuralnetGeoShapleyLeakLo}{57.6}
\newcommand{\RobNeuralnetGeoShapleyN}{5}
\newcommand{\RobNeuralnetGeoShapleyTopPen}{99.2}
\newcommand{\RobNeuralnetGeoShapleyTopPenHi}{101.4}
\newcommand{\RobNeuralnetGeoShapleyTopPenLo}{97.1}
\newcommand{\RobNeuralnetKernelSHAPKept}{0.67}
\newcommand{\RobNeuralnetKernelSHAPKeptHi}{0.74}
\newcommand{\RobNeuralnetKernelSHAPKeptLo}{0.60}
\newcommand{\RobNeuralnetKernelSHAPLeak}{53.6}
\newcommand{\RobNeuralnetKernelSHAPLeakHi}{56.7}
\newcommand{\RobNeuralnetKernelSHAPLeakLo}{50.6}
\newcommand{\RobNeuralnetKernelSHAPN}{5}
\newcommand{\RobNeuralnetKernelSHAPTopPen}{98.7}
\newcommand{\RobNeuralnetKernelSHAPTopPenHi}{101.0}
\newcommand{\RobNeuralnetKernelSHAPTopPenLo}{96.5}
\newcommand{\RobNeuralnetLIMEKept}{0.02}
\newcommand{\RobNeuralnetLIMEKeptHi}{0.03}
\newcommand{\RobNeuralnetLIMEKeptLo}{0.02}
\newcommand{\RobNeuralnetLIMELeak}{0.0}
\newcommand{\RobNeuralnetLIMELeakHi}{0.0}
\newcommand{\RobNeuralnetLIMELeakLo}{0.0}
\newcommand{\RobNeuralnetLIMEN}{5}
\newcommand{\RobNeuralnetLIMETopPen}{0.7}
\newcommand{\RobNeuralnetLIMETopPenHi}{2.6}
\newcommand{\RobNeuralnetLIMETopPenLo}{$-$1.2}
\newcommand{\RobPending}{none}
\newcommand{\RobRandomforestGeoShapleyKept}{0.72}
\newcommand{\RobRandomforestGeoShapleyKeptHi}{0.78}
\newcommand{\RobRandomforestGeoShapleyKeptLo}{0.65}
\newcommand{\RobRandomforestGeoShapleyLeak}{60.8}
\newcommand{\RobRandomforestGeoShapleyLeakHi}{64.1}
\newcommand{\RobRandomforestGeoShapleyLeakLo}{57.6}
\newcommand{\RobRandomforestGeoShapleyN}{5}
\newcommand{\RobRandomforestGeoShapleyTopPen}{100.0}
\newcommand{\RobRandomforestGeoShapleyTopPenHi}{100.0}
\newcommand{\RobRandomforestGeoShapleyTopPenLo}{100.0}
\newcommand{\RobRandomforestKernelSHAPKept}{0.67}
\newcommand{\RobRandomforestKernelSHAPKeptHi}{0.74}
\newcommand{\RobRandomforestKernelSHAPKeptLo}{0.60}
\newcommand{\RobRandomforestKernelSHAPLeak}{53.6}
\newcommand{\RobRandomforestKernelSHAPLeakHi}{56.7}
\newcommand{\RobRandomforestKernelSHAPLeakLo}{50.6}
\newcommand{\RobRandomforestKernelSHAPN}{5}
\newcommand{\RobRandomforestKernelSHAPTopPen}{100.0}
\newcommand{\RobRandomforestKernelSHAPTopPenHi}{100.0}
\newcommand{\RobRandomforestKernelSHAPTopPenLo}{100.0}
\newcommand{\RobRandomforestLIMEKept}{0.02}
\newcommand{\RobRandomforestLIMEKeptHi}{0.03}
\newcommand{\RobRandomforestLIMEKeptLo}{0.01}
\newcommand{\RobRandomforestLIMELeak}{0.0}
\newcommand{\RobRandomforestLIMELeakHi}{0.0}
\newcommand{\RobRandomforestLIMELeakLo}{0.0}
\newcommand{\RobRandomforestLIMEN}{5}
\newcommand{\RobRandomforestLIMETopPen}{0.7}
\newcommand{\RobRandomforestLIMETopPenHi}{2.6}
\newcommand{\RobRandomforestLIMETopPenLo}{$-$1.2}
\newcommand{\RobkmeansbackgroundGeoShapleyKept}{0.68}
\newcommand{\RobkmeansbackgroundGeoShapleyKeptHi}{0.75}
\newcommand{\RobkmeansbackgroundGeoShapleyKeptLo}{0.61}
\newcommand{\RobkmeansbackgroundGeoShapleyLeak}{59.5}
\newcommand{\RobkmeansbackgroundGeoShapleyLeakHi}{64.1}
\newcommand{\RobkmeansbackgroundGeoShapleyLeakLo}{54.9}
\newcommand{\RobkmeansbackgroundGeoShapleyN}{5}
\newcommand{\RobkmeansbackgroundGeoShapleyTopPen}{99.4}
\newcommand{\RobkmeansbackgroundGeoShapleyTopPenHi}{101.1}
\newcommand{\RobkmeansbackgroundGeoShapleyTopPenLo}{97.6}
\newcommand{\RobkmeansbackgroundKernelSHAPKept}{0.64}
\newcommand{\RobkmeansbackgroundKernelSHAPKeptHi}{0.71}
\newcommand{\RobkmeansbackgroundKernelSHAPKeptLo}{0.57}
\newcommand{\RobkmeansbackgroundKernelSHAPLeak}{52.0}
\newcommand{\RobkmeansbackgroundKernelSHAPLeakHi}{56.4}
\newcommand{\RobkmeansbackgroundKernelSHAPLeakLo}{47.5}
\newcommand{\RobkmeansbackgroundKernelSHAPN}{5}
\newcommand{\RobkmeansbackgroundKernelSHAPTopPen}{96.8}
\newcommand{\RobkmeansbackgroundKernelSHAPTopPenHi}{101.7}
\newcommand{\RobkmeansbackgroundKernelSHAPTopPenLo}{91.8}
\newcommand{\RuleCfiveFeatureGeoShapley}{leak depends on Moran's I}
\newcommand{\RuleCfiveFeatureKernelSHAP}{leak depends on Moran's I}
\newcommand{\RuleCfiveFeatureLIME}{no detectable dependence}
\newcommand{\RuleConeFeatureGeoShapley}{fooled}
\newcommand{\RuleConeFeatureKernelSHAP}{fooled}
\newcommand{\RuleConeFeatureLIME}{fooled}
\newcommand{\RuleConeLocationGeoShapley}{fooled}
\newcommand{\RuleConeLocationKernelSHAP}{fooled}
\newcommand{\RuleConeLocationLIME}{fooled}
\newcommand{\RuleConePrimeFeatureGeoShapley}{partly fooled}
\newcommand{\RuleConePrimeFeatureKernelSHAP}{partly fooled}
\newcommand{\RuleConePrimeFeatureLIME}{fully fooled}
\newcommand{\RuleConePrimeLocationGeoShapley}{partly fooled}
\newcommand{\RuleConePrimeLocationKernelSHAP}{partly fooled}
\newcommand{\RuleConePrimeLocationLIME}{fully fooled}
\newcommand{\RuleCtwoFeatureGeoMinusKs}{GeoShapley slightly harder to fool}
\newcommand{\RuleCtwoLocationGeoMinusKs}{no detectable difference at this sample size}
\newcommand{\RuleEsevenFeature}{go: include section}
\newcommand{\RuleMechFeatureGeoShapley}{prediction holds}
\newcommand{\RuleMechFeatureKernelSHAP}{prediction holds}
\newcommand{\RuleMechLocationGeoShapley}{prediction holds}
\newcommand{\SeedsFeature}{5}
\newcommand{\SeedsLocation}{5}
\newcommand{\SharePenFeature}{15.9}
\newcommand{\SharePenFeatureHi}{17.4}
\newcommand{\SharePenFeatureLo}{14.4}
\newcommand{\SharePenLocation}{15.9}
\newcommand{\SharePenLocationHi}{17.3}
\newcommand{\SharePenLocationLo}{14.5}
\newcommand{\SlopeMoranGeoShapley}{$-$0.11}
\newcommand{\SlopeMoranGeoShapleyHi}{$-$0.10}
\newcommand{\SlopeMoranGeoShapleyLo}{$-$0.12}
\newcommand{\SlopeMoranKernelSHAP}{$-$0.14}
\newcommand{\SlopeMoranKernelSHAPHi}{$-$0.06}
\newcommand{\SlopeMoranKernelSHAPLo}{$-$0.21}
\newcommand{\SlopeMoranLIME}{0.00}
\newcommand{\SlopeMoranLIMEHi}{0.00}
\newcommand{\SlopeMoranLIMELo}{0.00}
\newcommand{\TriggerMoran}{0.83}
\newcommand{\VerMachine}{x86\_64}
\newcommand{\VerNumpy}{2.2.6}
\newcommand{\VerPython}{3.10.20}
\newcommand{\VerShap}{0.49.1}
\newcommand{\VerSklearn}{1.7.2}
\newcommand{\VerXgb}{3.2.0}

\IfFileExists{tables/numbers_missing.tex}{% generated: macros used in the text but not produced by this run
}{}
\newcommand{\tabinput}[1]{\IfFileExists{tables/#1.tex}{\input{tables/#1}}{\todo{table \detokenize{#1} not generated yet}}}
\newcommand{\figinput}[3]{\IfFileExists{figures/#1.pdf}{\begin{figure}[t]\centering\includegraphics[width=#2\linewidth]{figures/#1}\caption{#3}\label{fig:#1}\end{figure}}{\todo{figure \detokenize{#1} not generated yet}}}

\title{The estimator, not the geography:\\ scaffold attacks hide a racial price penalty from LIME, KernelSHAP and GeoShapley\\[6pt]
\large Draft v1 for discussion with Peng Luo}
\author{Robert Romani}
\date{\today\ --- every number in this draft is generated from the result files; see Appendix C}

\begin{document}
\maketitle

\begin{abstract}
Post-hoc explanations are increasingly used to audit machine-learning models for spatial bias, and GeoShapley was
designed to separate location effects from other features. We show that a scaffold attack can hide a racial price
penalty from LIME, KernelSHAP and GeoShapley. A home-valuation model for King County, Washington, halves predicted
prices in census tracts where Black and Latino residents exceed a quarter of the population; a detector trained on
each explainer's own queries serves an innocuous model to the explainer and the biased model to real homes. LIME then
shows \FLIMEMatKept\ of the penalty while \FLIMEMatFid\% of real homes still receive it. KernelSHAP and GeoShapley show
\FKernelSHAPMatKept\ and \FGeoShapleyMatKept\ of it at about 90\% fidelity, and the attacker can hide more only by
withholding the penalty from real homes. GeoShapley is no harder to fool than KernelSHAP, and under attack its
location effect shrinks while spurious feature-by-location interactions appear. The difference between explainers is
set by their estimators, not by their treatment of space: efficiency keeps the penalty inside a Shapley explanation,
and the share kept is close to the share of coalition hybrids the detector cannot tell from real homes. \todo{one
clause on generality once the robustness runs are in: model families and audit settings, rule R.} A residual check
that needs no explainer finds the penalised homes (AUC \EsevenAucLagScaffoldTen\ for a 21\% penalty).

\end{abstract}

% =====================================================================================
\section{Introduction}

Machine-learning models now set or inform home valuations, insurance prices and lending decisions, and in each of
these location is among the strongest predictors. Location is also a proxy for race and ethnicity, so a valuation
model can reproduce the effects of historical redlining without ever being given a demographic variable
\citeneeded{redlining and algorithmic valuation}. Regulators, lenders and researchers increasingly check such models
with post-hoc explanation methods, which attribute each prediction to the model's inputs.

Geographers have adapted these methods to spatial data. GeoShapley \citep{li2024geoshapley} treats the coordinate
pair as one player in a Shapley game and separates an intrinsic location effect from location-invariant feature
effects and feature-by-location interactions. For a fairness audit this is attractive: if a model penalises a
neighbourhood, the location effect should show it.

That reading assumes the model cooperates. \citet{slack2020fooling} showed that LIME and SHAP can be deceived by a
\emph{scaffold}: a model that recognises the synthetic inputs an explainer generates and answers them with an
innocuous model, while real inputs receive biased predictions. The attack has been shown only on non-spatial tabular
benchmarks. Whether an explainer built for spatial data inherits the weakness, resists it, or fails in its own way
has not been examined.

We examine it on a home-valuation model for King County, Washington. The model halves predicted prices in census
tracts where Black and Latino residents exceed a quarter of the population, and a scaffold hides the penalty behind
a detector trained on each explainer's own queries. Our central finding is that \emph{how much a scaffold can hide is
set by the explainer's estimator --- where it queries the model --- and not by whether the explainer knows about
geography.} LIME samples far from real homes and can be fooled almost for free. KernelSHAP and GeoShapley satisfy
efficiency and evaluate the explained home itself, together with many near-real hybrids, so the attacker can hide
the penalty from them only by also withholding it from real homes. We measure that stealth--fidelity trade-off and
explain it with an argument about Shapley coalitions that accounts for the measured trade-off.

We ask five questions.
\begin{enumerate}
  \item Can a scaffold hide a demographic price penalty from LIME, KernelSHAP and GeoShapley in a real housing market?
  \item What does hiding cost the attacker, in the share of real homes that still receive the bias?
  \item Does GeoShapley's treatment of location make it harder to fool than a non-spatial Shapley explainer, and what
        happens to its decomposition under attack?
  \item Why do the explainers differ?
  \item Do the answers survive other model families, auditors who change the explainer's settings, a smaller penalty
        and weaker spatial autocorrelation of the biased attribute?
\end{enumerate}
In brief: all three explainers can be fooled. LIME is fooled completely, keeping \FLIMEMatKept\ of the penalty while
the scaffold still charges \FLIMEMatFidPen\% of the penalised homes; an earlier report that LIME resists the attack in
a geographic setting came from a detector that did not match LIME's sampler. KernelSHAP and GeoShapley keep
\FKernelSHAPMatKept\ and \FGeoShapleyMatKept\ of the penalty when the attacker serves the biased model to about 90\%
of real homes (\FKernelSHAPMatFidPen\% and \FGeoShapleyMatFidPen\% of the penalised ones), and the share they keep
falls roughly in step with the share of homes the attacker is willing to spare. On the penalised homes themselves the
trigger stays among the three largest attributions; the scaffold shrinks it and pushes it down the explainer's global
importance ranking, but cannot take it out of those homes' explanations. Treating location as a joint player protects
GeoShapley little: in the feature design \FCtwoClause, in the location design \LCtwoClause, and under attack its
location effect shrinks while its feature-by-location interactions roughly double. The share a Shapley explainer keeps
is close to the average, over coalition sizes, of the share of its hybrids that the detector routes to the biased
model. These results hold with random forests and neural networks and when the auditor changes the Shapley
background; only an auditor who switches LIME to a sampler the attacker did not anticipate recovers part of the
penalty. Finally, a residual check that uses no
explainer separates the penalised homes from the rest.

% =====================================================================================
\section{Background}

\paragraph{Perturbation-based explanation, and where each method queries the model.}
All three explainers estimate attributions by evaluating the model at inputs that are not real observations; this
is the attack surface.
\begin{itemize}
  \item \textbf{LIME} \citep{ribeiro2016lime} fits a weighted linear surrogate around an instance. Its default tabular
        sampler discretises each feature into training quartiles and draws new rows bin by bin, so almost every
        query is a combination of feature values that no real home has. The instance itself is the first row.
  \item \textbf{KernelSHAP} \citep{lundberg2017shap} estimates Shapley values by evaluating coalitions: the features in
        a coalition keep the instance's values and the rest are taken from a background row. With $M$ features it
        needs $2^M - 2$ coalitions per background row (2,046 in our feature design, 1,022 in the location design),
        which it enumerates in full.
  \item \textbf{GeoShapley} \citep{li2024geoshapley} uses the same coalition construction but moves the two coordinates
        together as one player, so every hybrid sits at a real location. It reports an intrinsic location effect
        (GEO), location-invariant primary effects and feature-by-location interaction effects that together satisfy
        efficiency.
\end{itemize}

\paragraph{The scaffold attack.}
\citet{slack2020fooling} train a classifier $d$ to tell real inputs from an explainer's synthetic queries and deploy
\begin{equation}
  e(x) = \begin{cases} f(x) & d(x) \ge \tau \quad\text{(looks real: biased answer)}\\
                       \psi(x) & d(x) < \tau \quad\text{(looks synthetic: innocuous answer).}\end{cases}
  \label{eq:scaffold}
\end{equation}
Real users receive biased predictions; the explainer, querying mostly synthetic points, reports the innocuous model.
The attack is only as good as the detector's match to the explainer's sampler, which is why \citet{slack2020fooling}
train one detector per explainer.

\paragraph{Related work.}
Explanations can be manipulated by modifying the model \citeneeded{Heo et al. 2019; Dombrowski et al. 2019}, and
rationalised to look fair \citeneeded{A\"ivodji et al. 2019, fairwashing}. Defences detect or avoid off-manifold
queries \citep{carmichael2023unfooling} \citeneeded{on-manifold Shapley, Frye et al.}. Shapley values in geography
\citep{li2022shap} and GeoShapley \citep{li2024geoshapley} are increasingly used to read spatial effects out of
machine-learning models. \todo{position against the geospatial XAI and algorithmic-bias literature Peng wants cited;
every entry to be checked against the published version}

% =====================================================================================
\section{Study area and data}

\paragraph{Housing.} 3,000 house sales in King County, Washington, from the sample distributed with the GeoShapley
package (\texttt{seattle\_sample\_3k.csv}) \todo{confirm the sale years of the sample}. Each sale has eight structural
features (bathrooms, living area, lot area, grade, condition, waterfront, view, age), projected coordinates (UTM zone
10N, metres) and the base-10 logarithm of the sale price (median sale \$435{,}000).

\paragraph{Demographics.} American Community Survey 2011--2015 table B03002 at census-tract level for the 398 King
County tracts, joined to 2015 TIGER/Line tract polygons. A point-in-polygon join assigns every sale to one of 334
tracts and attaches the tract share of Black (non-Hispanic) and Hispanic or Latino residents, \pct. Figure~\ref{fig:F1_study_area}
maps the share and the sales.

\paragraph{Split.} For each seed $s$ the 3,000 sales are split 80/20 into 2,400 training and 600 test homes. The
first 200 test homes are the homes explained; the other 400 calibrate the attacker's threshold (Section~\ref{sec:thresholds}).

\paragraph{Spatial structure.} The trigger is strongly autocorrelated: Moran's $I$ of \pct\ on an
8-nearest-neighbour, row-standardised weights matrix is \TriggerMoran.

\figinput{F1_study_area}{0.72}{Study area. Tracts shaded by the share of Black and Latino residents; tracts above the
0.25 threshold outlined; sales shown as points, red where the biased model applies its penalty.}

% =====================================================================================
\section{Methods}

\subsection{Honest, biased and innocuous models}
The honest model $h$ is a gradient-boosted regressor (XGBoost, 500 trees, depth 6, learning rate 0.1) on the eight
structural features and the two coordinates. The biased model subtracts a fixed penalty $\delta$ for homes in tracts
where the trigger exceeds 0.25:
\begin{equation}
  f(x) = h(x) - \delta\,\1\{T(x) > 0.25\}.
\end{equation}
With the price on a base-10 log scale, $\delta = 0.30$ multiplies the predicted price by $10^{-0.3} \approx 0.50$:
a penalty of one half.\footnote{Earlier project documents gave this as a 26\% penalty, which assumes natural
logarithms; Appendix~D and the repository's CHANGELOG list every correction of this kind.} The innocuous model $\psi$ is a second XGBoost regressor (300 trees,
depth 5) on the eight structural features only, so any explanation of it is free of location and demographics.
Table~\ref{tab:data} gives the fit of both on the test homes.

\paragraph{Two designs.} In the \emph{feature design} the trigger $T(x)$ is the \pct\ column, which the scaffold sees
as an eleventh feature. In the \emph{location design} there is no demographic column: the model sees the ten real
features, and $T(x)$ is the tract share looked up from the home's coordinates inside the model (a 25\,m raster of
the real tract polygons). This is redlining in the literal sense, and it is the case GeoShapley's location effect
exists to expose. A third design drops the coordinates altogether (LIME and KernelSHAP only).

\tabinput{tab_data}

\subsection{The scaffold and its detectors}
The deployed model is \eqref{eq:scaffold} with a 100-tree random forest as $d$. We build $d$ two ways.
\begin{itemize}
  \item \textbf{Generic detector}: trained on the perturbation recipes of \citet{slack2020fooling}: real rows plus
        Gaussian noise for LIME, and real rows with features swapped in from ten $k$-means centroids for the Shapley
        explainers. Threshold $\tau = 0.5$.
  \item \textbf{Matched detector}: trained on the queries each explainer actually sends. For LIME the negatives are
        draws from \texttt{LimeTabularExplainer}'s own sampler around every training home. For KernelSHAP and
        GeoShapley we run the explainer on training homes with a background drawn differently from the auditor's,
        record a random subset of every row it submits, and train $d$ to separate them from real training rows. Rows
        identical to a real home (the instance itself) are labelled real.
\end{itemize}
Detectors never see the test homes.

\subsection{Thresholds and fidelity}\label{sec:thresholds}
An attacker must decide how many real homes it is willing to misroute to $\psi$. We set the matched threshold on the
400 calibration homes: $\tau^* = \min(0.5, q_{0.10})$, where $q_{0.10}$ is the 10th percentile of their detector
scores, so at least 90\% of unseen real homes receive $f$. \emph{Fidelity} is then measured on the 200 explained homes,
which did not set $\tau^*$. Because $f = h$ on unpenalised homes, the bias actually delivered is the fidelity on the
penalised explained homes, which we report alongside. To trace the frontier we lower fidelity on purpose with the same
detector, at thresholds giving 75\% and 50\% fidelity on the calibration homes.

\subsection{Explainers}
LIME with 500 samples and the default quartile sampler; KernelSHAP with a 100-row background sample, every coalition
enumerated and no feature selection (\texttt{l1\_reg=False}, so the values are exact Shapley values for that
background)\footnote{shap's default \texttt{l1\_reg="num\_features(10)"} keeps at most ten features per home even under
full enumeration, which with eleven features zeroes one of them. At seed 42 switching it off changed the trigger's
attribution by less than 0.001.}; GeoShapley with the same 100-row background and the coordinates as the joint player.
GeoShapley is exact in the location design (eight non-location features); with nine non-location features (feature
design) the package evaluates all coalitions of size at most three and their complements plus Shapley-weighted draws,
about 600 of 1,024, which we seed per run.

\subsection{What we measure}
For explained home $i$ with attribution vector $\phi(x_i)$, trigger set $J$ (the \pct\ column; both UTM columns for
LIME and KernelSHAP in the location design; the location player for GeoShapley in the location design) and penalised
set $P = \{i : T(x_i) > 0.25\}$. GeoShapley's attributions enter in the package's Shapley-equivalent form, in which
each feature-by-location interaction is split equally between the feature and the location player; the
location player therefore carries the intrinsic GEO effect plus half of every interaction, and
Table~\ref{tab:geo} reports the GEO effect and the interactions separately. The measures are:
\begin{align}
  \text{signed attribution} \quad & a = \frac{1}{|P|}\sum_{i \in P}\sum_{j \in J}\phi_j(x_i), \\
  \text{penalty share kept} \quad & K = a^{\text{deployed}} / a^{\text{biased}}, \label{eq:K}\\
  \text{net of honest (location design)} \quad & K_{\text{net}} = (a^{\text{deployed}} - a^{\text{honest}}) / (a^{\text{biased}} - a^{\text{honest}}),\\
  \text{trigger in top three} \quad & \tfrac{1}{|A|}\textstyle\sum_{i\in A} \1\{\text{rank of } \textstyle\sum_{j\in J}|\phi_j(x_i)| \le 3\},\\
  \text{fidelity} \quad & \tfrac{1}{|A|}\textstyle\sum_{i\in A} \1\{d(x_i) \ge \tau\},\\
  \text{queries reaching } f \quad & \text{share of the explainer's own queries with } d \ge \tau,\\
  \text{detector AUC} \quad & \Pr\{d(x_{\text{real}}) > d(x_{\text{query}})\} + \tfrac12\Pr\{=\}.
\end{align}
where $A$ is either all $N = 200$ explained homes or the penalised ones, $P$; we report both. In the location design
the trigger score of LIME and KernelSHAP is the sum over both coordinates, ranked against single features, so their
top-three shares there are not directly comparable with GeoShapley's single location player. The reference in
\eqref{eq:K} is the explanation of the biased model served directly, with no scaffold: that, not the honest model, is
what the attack has to hide. $K = 1$ means the penalty is fully visible, $K = 0$ fully hidden.
Every detector score of every query is histogrammed, so AUC and the leak at any fidelity are computed exactly
without re-running an explainer.

\subsection{Why Shapley explainers resist: efficiency and coalition sizes}\label{sec:mechanism}
For a Shapley explainer with background $B$ the attributions of a home sum to the served prediction there minus its
background mean (efficiency):
\begin{equation}
  \sum_{j=1}^{M}\phi_j(x) = e(x) - \frac{1}{|B|}\sum_{b\in B} e(b).
\end{equation}
While the scaffold keeps fidelity, the explained home itself is routed to $f$ and the penalty sits in $e(x)$: it
cannot vanish from the explanation; it can only move to other features, or the attacker must stop serving $f$ at $x$.
LIME has no such constraint. The trigger's attribution is an average over coalition sizes with equal weight $1/M$,
\begin{equation}
  \phi_j(x) = \frac{1}{M}\sum_{k=0}^{M-1}\mathbb{E}_{|S|=k,\; j\notin S,\; b\sim B}\bigl[e(x_{S\cup j}, b) - e(x_S, b)\bigr],
  \label{eq:sizes}
\end{equation}
where $(x_S, b)$ takes the players in $S$ from $x$ and the rest from $b$. The largest coalition is the explained
home; near-complete coalitions differ from it in one or two players, and near-empty ones are a background (real)
home with one or two values swapped in. Both ends look real to a detector. If $\rho_k$ is the share of size-$k$
hybrids routed to $f$, the share of the penalty kept is roughly the average of $\rho_k$ over sizes, and the attacker
can lower it only by rejecting near-real points, which also rejects real homes. We estimate each term of
\eqref{eq:sizes} by Monte Carlo with the attack's own detector, background and thresholds (200 hybrids per home and
size, drawn independently of the attack run) and compare the result with the measured $K$ (rule M below). Because this
evaluates the same Shapley average with the same detector, agreement confirms the size decomposition and our
implementation rather than testing a separate hypothesis. The substantive approximation is the plain average over
sizes of $\rho_k$, which ignores how the penalty is spread across sizes; we report it as a description that was not
specified in advance.

\subsection{Robustness}
\begin{itemize}
  \item \textbf{Model family}: $h$ and $\psi$ as random forests (100 trees, leaf size 3) or neural networks (two hidden
        layers of 256 and 128 units, standardised inputs).
  \item \textbf{Audit settings the attacker did not train for}: Shapley background of 25 or 200 rows, or 100 $k$-means
        centroids; LIME with 5,000 samples or with its Gaussian sampler. The detector stays the one trained for the
        defaults; the biased reference uses the auditor's changed settings.
  \item \textbf{Spatial autocorrelation}: a share of homes receives the trigger value of a random other home, which
        keeps the marginal distribution and lowers Moran's $I$.
  \item \textbf{Smaller penalty}: $\delta = 0.10$ (a 21\% price cut).
\end{itemize}

\subsection{An outcome-based check}
No explainer is involved: on the 600 held-out homes we compute residuals $r = e(x) - y$ of the deployed model against
observed log price and ask whether they reveal the penalised homes, using the residual itself, its local mean over
the 8 nearest neighbours, global Moran's $I$, and local Moran (LISA) low--low clusters
\citep{moran1950,anselin1995lisa}, for penalties from 0.02 to 0.30. The scaffold here is the LIME-matched one, which
serves almost every real home with $f$.

\subsection{Statistics and pre-specified decision rules}
The seed is the unit of replication; a seed re-runs the split, the models, the detectors, the background draw and the
LIME sampling. Each metric is a mean over seeds with a 95\% interval $\bar x \pm t_{S-1}^{-1}(0.975)\,\hat\sigma/\sqrt S$;
contrasts are paired by seed. The rules below were fixed before the runs.
\begin{itemize}
  \item \textbf{C1 (fooled)}: $1 - K$ under the matched scaffold has an interval above 0, Holm-corrected over the three
        explainers.
  \item \textbf{C1$'$ (fully fooled)}: trigger in the top three for at most 5\% of the explained homes and $K \le
        0.10$ at fidelity at least 90\%; otherwise \emph{partly fooled}. (Computed on penalised homes only, the
        top-three part gives the same classes.)
  \item \textbf{C2}: GeoShapley minus KernelSHAP in $K$, matched scaffold; an interval above 0 reads as
        ``GeoShapley slightly harder to fool'', an interval containing 0 as no detectable difference at this sample
        size.
  \item \textbf{C3}: GeoShapley's GEO effect and interaction sums, biased against matched; a change is claimed only if
        the paired interval excludes 0.
  \item \textbf{C5}: slope of the matched leak on Moran's $I$ over at least three levels.
  \item \textbf{M}: the Monte Carlo prediction holds if its mean absolute error against measured $K$ is no larger than
        the seed-to-seed standard deviation of measured $K$; otherwise only the efficiency identity is claimed.
  \item \textbf{R}: a robustness configuration agrees with the main result if each explainer's C1$'$ class is unchanged.
  \item \textbf{E7}: the outcome-check section is included if, at $\delta = 0.10$, the scaffold's local-mean AUC is at
        least 0.80 and the honest model's LISA flag rate on unpenalised homes is at most 6\%.
\end{itemize}

% =====================================================================================
\section{Results}
% Every number below is a macro from tables/numbers.tex (scripts/make_paper_tables.py). Rule outcomes are quoted
% from the same file (\Rule... macros), so a changed outcome shows up in the text rather than contradicting it.

\subsection{Models and detectors}
The honest model explains the test prices well ($R^2$ \HonestRtwoFeature\ in the feature design, \HonestRtwoLocation\
in the location design) and the innocuous model, which sees neither location nor demographics, less well
(\PsiRtwoFeature\ and \PsiRtwoLocation; Table~\ref{tab:data}). The biased model penalises \SharePenFeature\% of test
homes, about \NPenFeature\ of the 200 explained homes in each seed. In the location design the model's coordinate
lookup puts \RasterAgree\% of test homes in the same tract as the spatial join in the data file.

The matched detectors separate LIME's queries from real homes completely (AUC \FLIMEMatAuc), but separate the
Shapley explainers' queries only partly (AUC \FKernelSHAPMatAuc\ for KernelSHAP and \FGeoShapleyMatAuc\ for
GeoShapley in the feature design). This difference, not the attack recipe, drives everything that follows.

\subsection{A scaffold hides the penalty from all three explainers}
\figinput{F2_main}{0.66}{Penalty share kept in the feature design (1 = the explanation shows the whole penalty, as it
does when the biased model is served directly; 0 = none of it). Mean and 95\% interval over seeds.}

\tabinput{tab_main_feature}

\paragraph{Feature design.} Served directly, the biased model's penalty is plain in every explanation: the trigger
is among the three most important features for \FLIMEBiasedTop\% (LIME), \FKernelSHAPBiasedTop\% (KernelSHAP) and
\FGeoShapleyBiasedTop\% (GeoShapley) of the explained homes, with a signed attribution on penalised homes of about
\FKernelSHAPBiasedSigned\ in log$_{10}$ price (Table~\ref{tab:main_feature}).

The generic detector of \citet{slack2020fooling} does not fool LIME here: it lets \FLIMEGenLeak\% of LIME's queries
through to $f$, and LIME keeps \FLIMEGenKept\ of the penalty. This is the source of the earlier project finding that
LIME resists the attack in a geographic setting (Appendix~D); it is a failure of the detector, not a property of
LIME. With a detector matched to LIME's own sampler, no LIME query reaches $f$ (\FLIMEMatLeak\%), the scaffold still
serves $f$ to \FLIMEMatFid\% of the explained homes, and LIME keeps \FLIMEMatKept\ of the penalty, with the trigger in
the top three for \FLIMEMatTop\% of homes. LIME is \RuleConePrimeFeatureLIME.

Against the Shapley explainers the matched scaffold, at \FKernelSHAPMatFid\% and \FGeoShapleyMatFid\% fidelity, cuts
the penalty to \FKernelSHAPMatKept\ (95\% interval \FKernelSHAPMatKeptLo--\FKernelSHAPMatKeptHi) for KernelSHAP and
\FGeoShapleyMatKept\ (\FGeoShapleyMatKeptLo--\FGeoShapleyMatKeptHi) for GeoShapley, and the trigger drops out of the
top three for about half of the homes (\FKernelSHAPMatTop\% and \FGeoShapleyMatTop\% remain). Rule C1 (Holm-corrected)
gives LIME \RuleConeFeatureLIME, KernelSHAP \RuleConeFeatureKernelSHAP\ and GeoShapley \RuleConeFeatureGeoShapley;
by rule C1$'$ KernelSHAP is \RuleConePrimeFeatureKernelSHAP\ and GeoShapley \RuleConePrimeFeatureGeoShapley. The detector cannot do better because \FKernelSHAPMatLeak\% of KernelSHAP's and
\FGeoShapleyMatLeak\% of GeoShapley's queries look real enough to be served by $f$.

\paragraph{Location design.} Here location has a legitimate price effect, so even the honest model's explanation
puts a negative location attribution on the homes the biased model penalises (\LKernelSHAPHonestSigned\ for
KernelSHAP). Table~\ref{tab:main_location} therefore gives the share kept both raw and net of the honest model. The
LIME-matched scaffold removes location from LIME's explanation altogether: the raw share is \LLIMEMatKeptRaw, location
is in the top three for \LLIMEMatTop\% of homes against \LLIMEHonestTop\% under the honest model, and the net share
is negative (\LLIMEMatKept): LIME now shows less location effect than an honest model would. KernelSHAP and
GeoShapley keep \LKernelSHAPMatKept\ and \LGeoShapleyMatKept\ of the penalty net of the honest model (raw
\LKernelSHAPMatKeptRaw\ and \LGeoShapleyMatKeptRaw), at fidelity \LKernelSHAPMatFid\% and \LGeoShapleyMatFid\%.
Location stays in their top three for most homes (\LKernelSHAPMatTop\% and \LGeoShapleyMatTop\%), as it should for a
housing model; what the scaffold removes is about half of the penalty carried by it. In this design rule C1 gives
LIME \RuleConeLocationLIME, KernelSHAP \RuleConeLocationKernelSHAP\ and GeoShapley \RuleConeLocationGeoShapley, and
rule C1$'$ gives LIME \RuleConePrimeLocationLIME, KernelSHAP \RuleConePrimeLocationKernelSHAP\ and GeoShapley
\RuleConePrimeLocationGeoShapley.

\tabinput{tab_main_location}

\subsection{What hiding costs: the stealth--fidelity frontier}
\figinput{F3_frontier}{0.92}{The stealth--fidelity frontier. Each filled marker is the matched detector at a stricter
threshold (90\%, 75\% and 50\% fidelity on the calibration homes); hollow markers are the generic detector at 0.5.
The location panel shows the share kept net of the honest model. Mean and 95\% interval over seeds in both
directions.}

An attacker who wants to hide more from a Shapley explainer has one lever: withhold the bias from more real homes.
Lowering the matched threshold so that 75\% and then 50\% of real homes receive $f$ lowers KernelSHAP's share kept
from \FKernelSHAPMatKept\ to \FKernelSHAPMatSeventyFiveKept\ and \FKernelSHAPMatFiftyKept, and GeoShapley's from
\FGeoShapleyMatKept\ to \FGeoShapleyMatSeventyFiveKept\ and \FGeoShapleyMatFiftyKept\ (Table~\ref{tab:frontier},
Figure~\ref{fig:F3_frontier}). The share kept falls roughly in step with fidelity. Even at 50\% fidelity --- a
scaffold that spares half the homes it was built to penalise --- the trigger is still in the top three for
\FKernelSHAPMatFiftyTop\% (KernelSHAP) and \FGeoShapleyMatFiftyTop\% (GeoShapley) of homes.

The generic detector, trained on different perturbations and run at a fixed threshold, reaches
\FKernelSHAPGenFid\% fidelity and a share kept of \FKernelSHAPGenKept\ (KernelSHAP) and \FGeoShapleyGenKept\
(GeoShapley), on the same curve as the matched detector within the seed intervals. Within the detectors we tried, the
trade-off is set by where the explainer queries the model, not by how the attacker builds its detector. In the
location design the net share falls to \LKernelSHAPMatSeventyFiveKept\ and \LGeoShapleyMatSeventyFiveKept\ at 75\%
fidelity and to about zero at 50\% (\LKernelSHAPMatFiftyKept, \LGeoShapleyMatFiftyKept).

\tabinput{tab_frontier}

\subsection{GeoShapley is fooled like KernelSHAP, and its decomposition changes}
\paragraph{No detectable difference from KernelSHAP (rule C2).} Paired by seed, GeoShapley minus KernelSHAP in the
share kept under the matched scaffold is \FGeoMinusKsKept\ (interval \FGeoMinusKsKeptLo\ to \FGeoMinusKsKeptHi) in the
feature design and \LGeoMinusKsKept\ (\LGeoMinusKsKeptLo\ to \LGeoMinusKsKeptHi) in the location design: feature
design, \RuleCtwoFeatureGeoMinusKs; location design, \RuleCtwoLocationGeoMinusKs. GeoShapley's queries are
somewhat harder to tell from real homes (difference in detector AUC \FGeoMinusKsAuc\ in the feature design), consistent
with its moving both coordinates together so that every hybrid sits at a real location, but this does not translate
into a detectable difference in what the attacker can hide. Treating location as a joint player does not make an
explainer harder to fool; the coalition construction it shares with KernelSHAP decides that.

\paragraph{The decomposition moves (rule C3).} Under the matched scaffold the intrinsic location effect on penalised
homes moves towards zero and the feature-by-location interactions grow (Table~\ref{tab:geo},
Figure~\ref{fig:F4_geoshapley_decomposition}). In the feature design GEO on penalised homes goes from
\FGeoBiasedGeoPen\ under the biased model to \FGeoMatGeoPen\ under the matched scaffold, and the per-home sum of
absolute interaction effects from \FGeoBiasedInt\ to \FGeoMatInt. In the location design GEO goes from
\LGeoBiasedGeoPen\ (honest model: \LGeoHonestGeoPen) to \LGeoMatGeoPen, and the interactions from \LGeoBiasedInt\ to
\LGeoMatInt. Each change has a paired interval that excludes zero. The scaffold's routing depends on how real a hybrid
looks rather than on any single feature, and GeoShapley reports the resulting discontinuities as features acting
differently at different locations. An auditor who reads only the GEO map sees less of the penalty than the location
player as a whole carries.

\figinput{F4_geoshapley_decomposition}{0.82}{GeoShapley's decomposition by deployment: the intrinsic location effect
(GEO) on penalised homes, and the per-home sum of absolute feature-by-location interaction effects. Mean and 95\%
interval over seeds.}

\tabinput{tab_geo}

\subsection{Why: efficiency and coalition sizes}
\paragraph{Efficiency.} Under the matched scaffold the Shapley explanations of penalised homes still carry
\FKernelSHAPMatGapTotal\ (KernelSHAP) and \FGeoShapleyMatGapTotal\ (GeoShapley) of the penalty gap the biased model's
explanation puts between penalised and unpenalised homes: \FKernelSHAPMatKept\ and \FGeoShapleyMatKept\ on the trigger
and \FKernelSHAPMatReloc\ and \FGeoShapleyMatReloc\ moved to other features, most of it to \FKernelSHAPMatRelocTo.
Because the background term in the efficiency identity is the same for every home, the rest can leave the
explanation only through explained homes that the scaffold itself serves with $\psi$. LIME's explanation carries
\FLIMEMatGapTotal\ of the gap. The penalty that efficiency keeps in a Shapley explanation can be moved but
not removed while the explained home is served by $f$.

\paragraph{Routing by coalition size.} Figure~\ref{fig:F7_mechanism} shows the share of KernelSHAP's hybrids routed
to $f$ at each coalition size. It is U-shaped: \MechFeatureKernelSHAPRhoFirst\ for a background home with only the
trigger swapped in, \MechFeatureKernelSHAPRhoLast\ for the explained home itself ($k = \MechFeatureKernelSHAPM$), and
\MechFeatureKernelSHAPRhoMin\ at its lowest ($k = \MechFeatureKernelSHAPRhoMinK$). Both ends of the coalition range
are near-real homes the detector cannot reject without rejecting real homes; only the middle is far enough from the
data to be routed away, and equation~\eqref{eq:sizes} gives every size the same weight. Stricter thresholds lower the
whole curve, the ends included, which is why hiding costs fidelity.

\paragraph{Prediction (rule M).} The independent Monte Carlo evaluation of equation~\eqref{eq:sizes} predicts the
measured share kept with a mean absolute error of \MechFeatureKernelSHAPMAE\ for KernelSHAP (seed-to-seed standard
deviation of the measured share \MechFeatureKernelSHAPSD; \RuleMechFeatureKernelSHAP), \MechFeatureGeoShapleyMAE\ for
GeoShapley in the feature design (\RuleMechFeatureGeoShapley) and \MechLocationGeoShapleyMAE\ in the location design
(\RuleMechLocationGeoShapley; Table~\ref{tab:mechanism}). The plain average of the routing rates over sizes, which
ignores how the penalty is spread across sizes, is within \MechFeatureKernelSHAPRhoMAE\ of the measured share for
KernelSHAP: to a first approximation, \emph{the share of the penalty a Shapley explainer keeps is the average share of
its hybrids the detector cannot reject}.

\figinput{F7_mechanism}{0.92}{Mechanism. Left: share of with-trigger hybrids routed to $f$ by coalition size, at the
three attacker thresholds (feature design). Right: predicted against measured penalty share kept, one point per seed
and threshold; crosses are the plain average of the routing rates.}

\tabinput{tab_mechanism}

\subsection{Robustness}
\paragraph{Model families and audit settings.} Table~\ref{tab:robust} and Figure~\ref{fig:F8_robustness} repeat the
feature-design attack with random forests (honest $R^2$ \HonestRtwoForest) and neural networks (\HonestRtwoNet) in
place of XGBoost, and against auditors who change the explainer's settings after the attacker has trained its
detector. In \RobHold\ of \RobCells\ explainer--configuration cells the C1$'$ class is the same as in the main result
(rule R); exceptions: \RobBreakList. \todo{one or two sentences on the direction of the settings effects once all
seeds are in: does a larger background or $k$-means background raise the share kept, does LIME's Gaussian sampler
escape the quartile-trained detector?}

\tabinput{tab_robust}
\figinput{F8_robustness}{0.85}{Penalty share kept under the matched scaffold across model families and audit
settings the attacker did not train for (feature design). Mean and 95\% interval over seeds.}

\paragraph{Spatial autocorrelation of the trigger.} \todo{C5 result from Table~\ref{tab:autocorr} and
Figure~\ref{fig:F6_autocorrelation}: slope of the matched leak on Moran's $I$ (\RuleCfiveFeature).}
\tabinput{tab_autocorr}
\figinput{F6_autocorrelation}{0.66}{Queries reaching $f$ under the matched scaffold as the trigger's spatial
autocorrelation is lowered. Mean and 95\% interval over seeds.}

\paragraph{A smaller penalty.} \todo{Table~\ref{tab:penalty010}: share kept at $\delta = 0.10$ (LIME \PenTenLIMEKept,
KernelSHAP \PenTenKernelSHAPKept, GeoShapley \PenTenGeoShapleyKept).}
\tabinput{tab_penalty010}

\subsection{An outcome-based check sees what the explainers miss}
No explainer is needed to find this penalty. On the 600 held-out homes, the 8-nearest-neighbour mean of the
deployed model's residuals separates penalised from unpenalised homes with AUC \EsevenScaffoldAucLagThirty\ at
$\delta = 0.30$ and \EsevenAucLagScaffoldTen\ (interval \EsevenAucLagScaffoldTenLo--\EsevenAucLagScaffoldTenHi) at
$\delta = 0.10$, for the scaffold that fools LIME completely (Table~\ref{tab:e7}, Figure~\ref{fig:F5_outcome_check}).
The check loses power only for penalties of a few percent (AUC \EsevenScaffoldAucLagTwo\ at $\delta = 0.02$). It costs
the auditor few false alarms: LISA flags \EsevenHonestFlag\% of the honest model's unpenalised homes as low--low
clusters. The scaffold cannot evade this check without changing what real homes are charged, because it is computed
from served predictions and observed prices only. By rule E7 (\RuleEsevenFeature) the check is reported here; it is a
complement to explanation-based audits, not a defence we evaluate against adaptive attackers.

\figinput{F5_outcome_check}{0.6}{Outcome-based check: AUC of the local mean residual for the penalised homes, by
penalty size, for the biased model and the LIME-matched scaffold; the line marks the honest model. Mean and 95\%
interval over seeds.}

\tabinput{tab_e7}


% =====================================================================================
\section{Discussion}
% Numbers are macros from tables/numbers.tex. Statements that depend on a rule outcome name the rule.

\paragraph{The estimator, not the geography.} What separates the explainers in our results is where each one queries
the model. LIME's quartile sampler produces rows that no real home resembles; a detector separates them perfectly,
and LIME then explains the innocuous model while the deployed model keeps penalising almost every real home. The
Shapley explainers are harder to fool for two reasons that have nothing to do with space. Efficiency ties their
attributions to the served prediction at the explained home, so the penalty cannot leave the explanation while that
home is served by the biased model; and their coalition hybrids include, at both ends of the size range, points that
differ from a real home in one or two values and that a detector cannot reject without rejecting real homes.
GeoShapley moves the two coordinates together, which keeps every hybrid at a real location and makes its queries
slightly harder to detect, but its other players are still recombined across homes, and the middle coalitions are as
far from the data as KernelSHAP's. Under rule C2 we find no detectable difference between the two in what the
attacker can hide. Treating location carefully is valuable for interpretation; it is not, on its own, a defence. The
question to ask of any explainer used in an audit is the one the routing curves answer: what share of its queries
looks like real data?

\paragraph{What partial hiding means for an audit.} ``Partly fooled'' is not ``safe''. Against the matched scaffold
KernelSHAP and GeoShapley keep between half and three quarters of the penalty, and an auditor who knows what to look
for would still find the trigger among the leading features for about half of the penalised homes. But an audit is
rarely a search for one named variable, and an attacker who accepts lower fidelity can do much better: at 50\%
fidelity the scaffold still charges half of the targeted homes the full penalty while KernelSHAP shows
\FKernelSHAPMatFiftyKept\ of it. The frontier in Figure~\ref{fig:F3_frontier} turns this into a quantity a regulator
can reason about: for a given explainer, how much bias can be hidden at what cost in bias actually delivered.

\paragraph{GeoShapley under attack.} For geographers the more specific lesson is about how GeoShapley fails. Under
the scaffold the intrinsic location effect on penalised homes shrinks towards zero and the feature-by-location
interactions roughly double (rule C3). A GEO map read on its own would understate the penalty, and the interaction
terms would invite a substantive story --- features that matter differently in different places --- that is in fact
an artefact of the attacker's routing. We did not set out to build a detector and do not evaluate one, but inflated
interactions without a corresponding change in the model's fit are a pattern worth examining in future work.

\paragraph{Reconciling with earlier results.} \citet{slack2020fooling} found LIME easier to fool than SHAP on
tabular benchmarks; we find the same ordering in a housing market. An earlier version of this project reported the
opposite, that LIME resists the attack in a geographic setting. That result came from the generic detector, which is
trained on Gaussian noise and lets \FLIMEGenLeak\% of LIME's real queries through. Replications of scaffold attacks
should match the detector to the explainer's actual sampler, or report how many queries reach the biased model.

\paragraph{Changing the audit settings.} \todo{after the robustness runs: which audit-setting changes, if any, move
the share kept (Table~\ref{tab:robust}); whether an auditor gains anything from a larger or $k$-means background or
from LIME's Gaussian sampler; keep the claim to what rule R and the intervals support.}

\paragraph{Looking at outcomes.} The penalty that the scaffold hides from explainers is plainly visible in the
residuals of the model's own predictions: their local mean separates penalised homes with AUC
\EsevenAucLagScaffoldTen\ even at a 21\% penalty, using tools --- local means, Moran's $I$, LISA --- that are standard
in geography. This check needs observed outcomes, loses power for penalties of a few percent, and we have not tested
it against an attacker who anticipates it. It complements explanation-based audits rather than replacing them, and it
reflects a broader point: for spatial bias, an audit that looks at what a model does to places can be more robust than
one that asks the model to explain itself.


\section{Limitations}
\begin{itemize}
  \item One metropolitan area in this draft. \todo{Cook County runs for the submission version}
  \item The penalty is synthetic and large; Table~\ref{tab:penalty010} repeats the main result at a smaller one.
  \item The attacker knows which explainer is used; Table~\ref{tab:robust} varies the settings it does not know.
  \item The adversary is static: defences and attackers that adapt to them are out of scope.
  \item The auditor's Shapley background is drawn from the training homes, which the matched detector saw as real
        during training; near-empty coalitions may look real partly because the detector has memorised those rows.
        A $k$-means background, whose rows no real home matches exactly, changes the share kept little
        (Table~\ref{tab:robust}), but a background from homes the attacker never saw would be the cleaner test.
  \item 200 explained homes per seed, of which about \NPenFeature\ are penalised; the number of seeds is in every table.
  \item Feature-design GeoShapley uses the package's sampled coalitions, so its values carry Monte Carlo error that the
        seed intervals include.
\end{itemize}

\section{Conclusion}
A scaffold can hide a racial price penalty from the three explainers most used in geographic machine learning. How
much it can hide is set by each explainer's estimator: LIME, whose queries are unlike real homes, is fooled almost for
free, while KernelSHAP and GeoShapley, bound by efficiency and querying many near-real hybrids, can be fooled only at a
cost in the bias actually delivered. Making an explainer spatially explicit does not change this. Under attack
GeoShapley's location effect shrinks and spurious interactions appear, so its decomposition should not be read as
evidence that a model treats places fairly. Audits of spatial models should ask where an explainer queries the model,
and should pair explanations with checks on outcomes that a model cannot route around.


\section*{Data and code availability}
Housing sample: \texttt{seattle\_sample\_3k.csv} from the GeoShapley repository. Demographics: ACS 2011--2015 table
B03002 and 2015 TIGER/Line tracts, US Census Bureau. Code: the scripts listed in Appendix~C, run with fixed seeds;
repository \texttt{geospatial-xai-attacks}, branch \texttt{paper-draft-v1}.

\bibliographystyle{plainnat}
\begin{thebibliography}{99}
\bibitem[Anselin(1995)]{anselin1995lisa} Anselin, L. (1995). Local indicators of spatial association---LISA.
  \emph{Geographical Analysis} 27(2), 93--115. \todo{check}
\bibitem[Carmichael and Scheirer(2023)]{carmichael2023unfooling} Carmichael, Z. and Scheirer, W. J. (2023).
  Unfooling perturbation-based post hoc explainers. \emph{Proceedings of the AAAI Conference on Artificial
  Intelligence} 37(6).
\bibitem[Li(2022)]{li2022shap} Li, Z. (2022). Extracting spatial effects from machine learning model using local
  interpretation method: An example of SHAP and XGBoost. \emph{Computers, Environment and Urban Systems} 96, 101845.
\bibitem[Li(2024)]{li2024geoshapley} Li, Z. (2024). GeoShapley: A game theory approach to measuring spatial effects
  in machine learning models. \emph{Annals of the American Association of Geographers} 114(7), 1365--1385.
\bibitem[Lundberg and Lee(2017)]{lundberg2017shap} Lundberg, S. M. and Lee, S.-I. (2017). A unified approach to
  interpreting model predictions. \emph{Advances in Neural Information Processing Systems} 30.
\bibitem[Moran(1950)]{moran1950} Moran, P. A. P. (1950). Notes on continuous stochastic phenomena.
  \emph{Biometrika} 37(1/2), 17--23. \todo{check}
\bibitem[Ribeiro et al.(2016)]{ribeiro2016lime} Ribeiro, M. T., Singh, S. and Guestrin, C. (2016). ``Why should I
  trust you?'': Explaining the predictions of any classifier. \emph{Proceedings of KDD}, 1135--1144.
\bibitem[Slack et al.(2020)]{slack2020fooling} Slack, D., Hilgard, S., Jia, E., Singh, S. and Lakkaraju, H. (2020).
  Fooling LIME and SHAP: Adversarial attacks on post hoc explanation methods. \emph{Proceedings of the AAAI/ACM
  Conference on AI, Ethics, and Society}, 180--186.
\end{thebibliography}

\appendix
\section{Per-seed results}
\tabinput{tab_perseed}

\section{Detector diagnostics}
\tabinput{tab_detectors}
\tabinput{tab_main_nocoord}

\section{Reproducibility}
\paragraph{Code.} All in \texttt{scripts/}:
\begin{itemize}
  \item \texttt{run\_matched\_attack.py}: one design and seed per call; resumable; every detector and every block of
        25 explained homes is cached; all writes are atomic;
  \item \texttt{summarize\_matched\_attack.py}: seed-level tables;
  \item \texttt{mechanism\_routing.py}: Section~\ref{sec:mechanism};
  \item \texttt{run\_residual\_diagnostic.py}: the outcome check;
  \item \texttt{make\_paper\_tables.py} and \texttt{make\_paper\_figures\_matched.py}: every table, figure and
        number in this draft.
\end{itemize}
The job lists that were run are in \texttt{docs/paper/data/queues\_run.txt}.

\paragraph{Environment.} Python \VerPython\ on \VerMachine, numpy \VerNumpy, scikit-learn \VerSklearn, XGBoost
\VerXgb, shap \VerShap, lime 0.2.0.1, geoshapley 0.2.1.1. Before any run we reproduced the reference run at seed 42,
made earlier on a Linux aarch64 machine with the same package versions: honest $R^2$, innocuous $R^2$, LIME query
routing under both detectors, and KernelSHAP's trigger attribution for the biased model and the matched scaffold
agreed to at least twelve decimal places, and the matched KernelSHAP detector routed exactly the same share of
queries. Mean compute per seed was about \MinutesPerSeed\ minutes of explainer time on one core.

\paragraph{Changes during the audit of 12 November.} (i) KernelSHAP is run with \texttt{l1\_reg=False}; (ii)
GeoShapley's coalition sampling is seeded; (iii) all seeds were recomputed from scratch with the final code, so no
table mixes code versions or machines.

\section{Reconciliation with earlier project results}
\begin{itemize}
  \item \textbf{``LIME resists the attack in a geographic setting.''} Withdrawn. The generic (Slack-recipe) LIME detector
        is trained on Gaussian noise in raw units and lets \FLIMEGenLeak\% of LIME's real queries through to $f$
        (Table~\ref{tab:main_feature}); LIME then simply explains the biased model. With a detector matched to LIME's
        sampler, LIME is the most completely fooled of the three explainers, the ordering \citet{slack2020fooling}
        report.
  \item \textbf{``GeoShapley location attribution falls by 93\%.''} Withdrawn. It was computed with a retired custom
        GeoSHAP and divided an adversarial GeoSHAP value by an honest KernelSHAP one; the October summary already
        corrected it to $-58\%$ under the generic detector.
  \item \textbf{``A 26\% price penalty.''} Corrected: prices are on a base-10 log scale, so the 0.30 penalty halves the
        predicted price and the 0.10 penalty cuts it by 21\%.
  \item \textbf{16 October attack-settings table.} Its ``SHAP-matched'' detector reached high hiding only by routing
        7.5\% of real homes to $f$; the threshold rule here fixes fidelity at 90\% or more before hiding is measured.
  \item \textbf{Draft v0 (synthetic stand-in for the trigger, one seed).} Its KernelSHAP and GeoShapley figures
        (trigger in the top three for 13\% and 16\% of homes under the matched scaffold) came from a synthetic trigger
        column and a single seed; on the real column the matched scaffold leaves the trigger in the top three for
        \FKernelSHAPMatTop\% and \FGeoShapleyMatTop\% of homes (Table~\ref{tab:main_feature}).
\end{itemize}


\end{document}
